\chapter{高等线性代数与矩阵论}
\label{ch:linalg}

% ════════════════════════════════════════════════════════════════
%  数学基础部：为 SLAM/估计/规划/控制各部打线性代数地基。
%  自包含融会贯通：向量空间·线性变换·对偶·内积与伴随·谱定理·SVD·
%  极分解·Jordan 标准形·张量积与外代数。为 ch:lie / ch:rigid_body /
%  P1 优化（法方程/Cholesky）铺底。记号与 lie_theory/rigid_body 严格一致。
% ════════════════════════════════════════════════════════════════
% 本章局部算子宏（\providecommand 仅在未定义时生效，不覆盖全局/preamble）。
\providecommand{\im}{\operatorname{im}}
\providecommand{\Hom}{\operatorname{Hom}}
\providecommand{\tr}{\operatorname{tr}}
\providecommand{\diag}{\operatorname{diag}}
\providecommand{\proj}{\operatorname{proj}}
\providecommand{\Exp}{\operatorname{Exp}}
\providecommand{\Log}{\operatorname{Log}}

\begin{practice}[前置自测]
开始之前，先试答下列问题。若能流畅作答，可只速读本章表格与速查卡；若卡住，正是本章要为你解决的。
\begin{enumerate}
  \item 旋转矩阵 $\mathbf{R}\in\SO(3)$ 的“逆恰好等于转置”，这与“正交矩阵保内积”有什么关系？为什么“保内积”能推出“逆等于转置”？
  \item 机械臂静力学公式 $\boldsymbol{\tau}=\mathbf{J}^\top\mathbf{F}$（关节力矩 = 雅可比转置乘末端力）里的 $\mathbf{J}^\top$，是“矩阵转置”、“对偶映射”、还是“内积伴随”？三者何时一致、何时不同？
  \item 给定含噪的近似旋转矩阵 $\tilde{\mathbf{R}}$（$\tilde{\mathbf{R}}^\top\tilde{\mathbf{R}}\neq\mathbf{I}$），如何把它“投影回”离它最近的合法旋转？为什么逐列归一化或 Gram--Schmidt 都不是最优解？
  \item 罗德里格斯公式 $\mathbf{R}=\mathbf{I}+\sin\theta\,\mathbf{a}^\wedge+(1-\cos\theta)(\mathbf{a}^\wedge)^2$ 为什么只有三项、而非无穷级数？这与 $(\mathbf{a}^\wedge)^3=-\mathbf{a}^\wedge$ 有何关系？
  \item 最小二乘问题 $\min_{\mathbf{x}}\|\mathbf{A}\mathbf{x}-\mathbf{b}\|^2$ 的“法方程” $\mathbf{A}^\top\mathbf{A}\mathbf{x}=\mathbf{A}^\top\mathbf{b}$ 从何而来？它与“把 $\mathbf{b}$ 正交投影到 $\mathbf{A}$ 的列空间”是同一件事吗？
\end{enumerate}
\end{practice}

\section*{本章目标}
学完本章，你应当能够：
\begin{enumerate}
  \item \textbf{以坐标无关的语言}重建线性代数：向量空间、子空间、商空间、线性映射的核--像分解与秩--零度定理，并理解“矩阵只是基下的影子”；
  \item \textbf{构造并使用对偶空间} $V^*$、对偶映射 $\mathbf{T}^\top$、零化子 $W^\circ$，说清 twist/wrench 对偶、功率配对 $P=\mathbf{F}^\top\mathbf{V}$ 与 $\mathbf{J}^\top$ 的几何本质；
  \item \textbf{在内积空间上}做正交分解、Gram--Schmidt/QR、正交投影，由 Riesz 表示导出\emph{伴随算子} $\mathbf{A}^*$，并把它与对偶映射严格区分；
  \item \textbf{陈述并证明谱定理}（实对称、复正规、Schur 三角化），掌握正定性的特征值判据与对称半正定平方根；
  \item \textbf{从 $\mathbf{A}^\top\mathbf{A}$ 的谱分解构造 SVD}，理解“任意线性映射 = 旋转 $\times$ 各向异性拉伸 $\times$ 旋转”，并据此读出四个基本子空间、可操作度椭球、条件数；
  \item \textbf{由 SVD 导出极分解} $\mathbf{A}=\mathbf{Q}\mathbf{P}$、Moore--Penrose 伪逆 $\mathbf{A}^+$、正交 Procrustes/Kabsch 解，落地到旋转重正交化与点云配准；
  \item \textbf{掌握极小多项式、Cayley--Hamilton 与 Jordan 标准形}，理解它们为何是 $e^{\mathbf{A}t}$、罗德里格斯闭式与线性系统稳定性的代数源头；
  \item \textbf{用张量积与外代数}给出行列式的坐标无关定义、楔积与 Hodge 星，解释 $\so(3)\cong\Lambda^2\mathbb{R}^3$、叉积、惯性张量与 Plücker 坐标的代数统一。
\end{enumerate}

\subsection*{本章知识导航}
本章是一条“\emph{把朴素的矩阵代数升级为坐标无关、结构优先、对偶自觉的算子论}”的主线。它从最抽象的向量空间出发，逐层加装结构（对偶 $\to$ 内积 $\to$ 谱 $\to$ 分解），最后用多线性代数收口。结构如下：

\begin{center}
\small
\begin{tikzpicture}[
  node distance=4.5mm and 7mm, >=Stealth,
  blk/.style={draw, rounded corners, align=center, font=\small, inner sep=3pt, fill=main!6, draw=main!60!black},
  grp/.style={draw, rounded corners, align=center, font=\small, inner sep=3pt, fill=second!6, draw=second!60!black},
  thr/.style={draw, rounded corners, align=center, font=\small, inner sep=3pt, fill=third!8, draw=third!60!black},
  ln/.style={->, thick, gray!60!black}]
\node[blk] (vs) {向量空间\\ 子空间·商·秩零度};
\node[blk, right=of vs] (dual) {对偶空间\\ $V^*,\ \mathbf{T}^\top,\ W^\circ$};
\node[grp, below=7mm of vs] (inner) {内积空间\\ 正交·投影·QR};
\node[grp, right=of inner] (adj) {伴随算子\\ $\mathbf{A}^*$（Riesz）};
\node[thr, below=7mm of inner] (spec) {谱定理\\ 对称/正规/Schur};
\node[thr, right=of spec] (svd) {SVD·极分解\\ 伪逆·Procrustes};
\node[blk, below=7mm of spec, fill=main!10] (jordan) {极小多项式\\ Cayley--Hamilton·Jordan};
\node[blk, right=of jordan, fill=second!10] (tensor) {张量积·外代数\\ $\det$·Hodge·$\so(3)$};
\draw[ln] (vs)--(dual); \draw[ln] (vs.south)--(inner.north);
\draw[ln] (dual.south) -- (adj.north); \draw[ln] (inner)--(adj);
\draw[ln] (inner.south)--(spec.north); \draw[ln] (adj.south)--(svd.north);
\draw[ln] (spec)--(svd);
\draw[ln] (spec.south)--(jordan.north);
\draw[ln] (dual.east) -- ++(6mm,0) |- (tensor.east);
\draw[ln] (svd.south)--(tensor.north);
\end{tikzpicture}
\end{center}

\noindent\textbf{推荐路径}：第 1--2 节（向量空间、对偶）是\emph{所有}后续内容的语义底座，务必精读；第 3--4 节（内积、伴随）把代数空间几何化，是状态估计与最小二乘的母体；第 5--6 节（谱定理、SVD/极分解）是机器人日常计算的\emph{核心工具箱}，SLAM 后端、ICP、可操作度分析反复引用；第 7 节（Jordan）解释一切矩阵指数闭式的代数根源，对接 \cref{ch:lie} 的罗德里格斯公式；第 8 节（张量/外代数）给出行列式与 $\so(3)$ 的代数统一，为 \cref{ch:lie}、微分几何铺路。带（选读）标记者可在需要时回看。

\subsection*{前置知识桥接}
本章只假定读者熟悉本科线性代数的\emph{计算}层面（解方程、矩阵乘法、行列式的排列公式），不假定抽象代数。我们要做的是把这些计算技巧\emph{重新理解}为坐标无关的结构。\cref{ch:rigid_body} 已用到旋转矩阵 $\mathbf{R}\in\SO(3)$、反对称矩阵 $\mathbf{a}^\wedge$、叉积；\cref{ch:lie} 将密集使用本章的内积/伴随/谱/Jordan 工具。换言之，本章是“为何机器人学的矩阵长这样”的统一回答。

\subsection*{如果跳过本章会怎样}
\begin{itemize}
  \item 在 \cref{ch:lie} 里你会看到罗德里格斯公式“恰好”收敛成三项闭式，却不知道这是 $\so(3)$ 极小多项式 $x^3+x=0$（单位轴时）截断了指数级数——只能死记而无法推广到新李群；
  \item 在 P1 非线性优化里你会写出法方程 $\mathbf{H}^\top\boldsymbol{\Sigma}^{-1}\mathbf{H}\,\delta\mathbf{x}=\dots$ 与 Cholesky 分解，却不理解它是“正交投影”与“对称正定平方根”的代数化身，调不动病态与奇异；
  \item 在点云配准（ICP）里你会照抄 $\mathbf{R}=\mathbf{V}\mathbf{U}^\top$ 与那个神秘的 $\det$ 修正项，机器人“镜像翻转”时无从调试。
\end{itemize}

\begin{table}[H]\centering\small
\caption{本章阅读方式建议}
\begin{tabular}{lll}
\toprule
阅读方式 & 时间 & 适合谁 \\
\midrule
精读（含全部推导与练习） & 8--10 小时 & 首次系统建立算子论视角、要做后端/估计推导 \\
速读（跳过冗长证明细节） & 3 小时 & 有线性代数基础、想对齐本书记号与结构观 \\
速查（只看小结的符号表 + 定理速查表） & 20 分钟 & 写代码/推公式时回来查 \\
\bottomrule
\end{tabular}
\end{table}

\begin{note}
\textbf{本章记号与约定.}\;
域记 $\mathbb{F}$（常取 $\mathbb{R}$，频域分析取 $\mathbb{C}$）。向量/矩阵粗体：$\mathbf{x},\mathbf{A}$；向量空间 $V,W,U$。
对偶空间 $V^*=\Hom(V,\mathbb{F})$，对偶（转置）映射 $\mathbf{T}^\top:W^*\to V^*$，零化子 $W^\circ$。
内积 $\langle\cdot,\cdot\rangle$（\textbf{第一变量线性、第二变量共轭线性}，与数学惯例及七大教材一致）；诱导范数 $\|\mathbf{v}\|=\sqrt{\langle\mathbf{v},\mathbf{v}\rangle}$。
伴随算子记 $\mathbf{A}^*$（实标准内积下 $\mathbf{A}^*=\mathbf{A}^\top$，复 $\mathbf{A}^*=\bar{\mathbf{A}}^\top$）。
旋转群 $\SO(3)$、位姿群 $\SE(3)$、李代数 $\so(3),\se(3)$，反对称化 $(\cdot)^\wedge$、其逆 $(\cdot)^\vee$，指数/对数 $\Exp,\Log$——\textbf{与 \cref{ch:rigid_body}、\cref{ch:lie} 严格一致}。
奇异值降序 $\sigma_1\ge\sigma_2\ge\cdots\ge0$；特征值 $\lambda_i$；对角阵 $\boldsymbol{\Lambda}=\diag(\lambda_i)$、$\boldsymbol{\Sigma}=\diag(\sigma_i)$。
\end{note}

% ════════════════════════════════════════════════════════════════
\section{向量空间、线性映射与商空间}
\label{sec:vector-space}

\paragraph{动机：机器人学里的“向量”远不止箭头。}
位形空间的切向量（关节速度）、李代数 $\so(3)$ 的旋转向量、控制信号、协方差所在的对称矩阵空间——这些“行为像向量”的对象既不是箭头也不是数组。要让线性代数的全部定理\emph{一次证明、处处适用}，必须把“向量”公理化：抽掉它的具体形状，只保留“可加、可数乘、且这两个运算彼此相容”这一结构骨架。

\begin{definition}[向量空间]\label{def:vector-space}
设 $\mathbb{F}$ 为域。$\mathbb{F}$-\emph{向量空间} 是集合 $V$ 配两个运算——向量加法 $+:V\times V\to V$ 与标量乘法 $\cdot:\mathbb{F}\times V\to V$——使得 $(V,+)$ 是交换群（结合、交换、有零元 $\mathbf{0}$、有逆元），且标量乘法满足相容律：$a(b\mathbf{v})=(ab)\mathbf{v}$、$1\mathbf{v}=\mathbf{v}$、$a(\mathbf{u}+\mathbf{v})=a\mathbf{u}+a\mathbf{v}$、$(a+b)\mathbf{v}=a\mathbf{v}+b\mathbf{v}$（共八条公理）。
\end{definition}

\noindent “标量取自\emph{域}”至关重要：域中每个非零元有乘法逆，才保证 $a\mathbf{v}=\mathbf{0}\ (a\neq0)\Rightarrow\mathbf{v}=\mathbf{0}$（可以“除以 $a$”）。若标量只取自环（如 $\mathbb{Z}$），得到的是\emph{模}，模未必有基，结构远更复杂。常见向量空间：$\mathbb{F}^n$、次数 $\le n$ 的多项式 $\mathbb{F}[x]_{\le n}$（基 $\{1,x,\dots,x^n\}$）、$m\times n$ 矩阵 $M_{m\times n}(\mathbb{F})$、连续函数 $\mathcal{C}([0,1])$（\emph{无限维}）。

\paragraph{子空间、和与直和。}
$V$ 的非空子集 $W$ 若对加法与数乘封闭（等价判据：$\mathbf{0}\in W$ 且 $\forall\mathbf{u},\mathbf{v}\in W,a\in\mathbb{F}$ 有 $a\mathbf{u}+\mathbf{v}\in W$）即为\emph{子空间}，记 $W\le V$。两子空间的\emph{交}仍是子空间，但\emph{并一般不是}（$x$ 轴并 $y$ 轴不含 $(1,1)$）。把“并”修正为子空间的方式是\emph{和}
\begin{equation}
  W_1+W_2:=\{\mathbf{w}_1+\mathbf{w}_2:\mathbf{w}_i\in W_i\},
  \label{eq:subspace-sum}
\end{equation}
它是包含 $W_1\cup W_2$ 的最小子空间。维数遵循\emph{容斥}式：
\begin{equation}
  \dim(W_1+W_2)=\dim W_1+\dim W_2-\dim(W_1\cap W_2).
  \label{eq:dim-formula}
\end{equation}
当 $W_1\cap W_2=\{\mathbf{0}\}$ 时（等价：每个 $\mathbf{v}$ 有\emph{唯一}分解 $\mathbf{v}=\mathbf{w}_1+\mathbf{w}_2$；等价：维数恰好相加），称\emph{内直和} $V=W_1\oplus W_2$。

\begin{insight}[补空间不唯一，正反映“裸空间无正交”]
给定 $W\le V$，使 $V=W\oplus U$ 的\emph{补空间} $U$ 总存在（把 $W$ 的基扩成 $V$ 的基，扩出来的部分张成一个补）但\emph{不唯一}：$\mathbb{R}^2$ 中 $x$ 轴的补可以是 $y$ 轴、也可以是直线 $y=x$。补空间唯一需要一个“最自然”的方向——\emph{正交}补，而正交需要内积（\cref{sec:inner-product}）。\textbf{补空间的不唯一性恰恰是“纯向量空间没有内积”的几何写照}。机器人学里这不是细节：冗余机械臂关节速度空间 $\mathbb{R}^n=\ker\mathbf{J}\oplus U$ 的不同补 $U$ 对应不同伪逆——“用哪个伪逆”本质就是“选哪个补空间”。
\end{insight}

\paragraph{商空间：把一个子空间“压成一点”。}
若想“忽略 $W$ 方向的信息、只看离 $W$ 多远”，就把相差一个 $W$ 中向量的点视为同一点。定义等价关系 $\mathbf{v}\sim\mathbf{u}\iff\mathbf{v}-\mathbf{u}\in W$，等价类 $\mathbf{v}+W$ 称\emph{陪集}，全体陪集构成\emph{商空间} $V/W$，运算 $(\mathbf{u}+W)+(\mathbf{v}+W):=(\mathbf{u}+\mathbf{v})+W$、$a(\mathbf{v}+W):=(a\mathbf{v})+W$。

\begin{proposition}[商空间维数与良定义]\label{prop:quotient}
$V/W$ 上的运算\emph{良定义}（不依赖陪集代表元的选取）：若 $\mathbf{u}_1-\mathbf{u}_2,\mathbf{v}_1-\mathbf{v}_2\in W$，则 $(\mathbf{u}_1+\mathbf{v}_1)-(\mathbf{u}_2+\mathbf{v}_2)\in W$。有限维时 $\dim(V/W)=\dim V-\dim W$。
\end{proposition}

\begin{pitfall}[概念误区：把 $V/W$ 当成 $V$ 的子空间]
$V/W$ 的元素是\emph{陪集}（$V$ 的子集），不是 $V$ 的向量；$V/W$ 是一个全新的向量空间，\emph{不住在} $V$ 里。它与任一补空间 $U$ 同构（$U\cong V/W$），但 $V/W$ 由 $V,W$ \textbf{唯一}确定、不需任何选择，而补空间不唯一。直觉：$V/W$ 记录“海拔”（相对 $W$ 的高度），补空间是“选一条垂直线挑代表点”。良定义性验证（运算不依赖代表元）是凡涉及“陪集 / 选代表元”的构造都\emph{必须}做的一步。
\end{pitfall}

\paragraph{线性映射、核与像。}
保持加法与数乘的映射 $\mathbf{T}:V\to W$ 称\emph{线性映射}。它由“基的像”唯一确定（\emph{基的泛性质}）。两个基本子空间：\emph{核} $\ker\mathbf{T}=\{\mathbf{v}:\mathbf{T}\mathbf{v}=\mathbf{0}\}\le V$（被压成零的部分）与\emph{像} $\im\mathbf{T}=\{\mathbf{T}\mathbf{v}\}\le W$（能到达的部分）。

\begin{theorem}[秩--零度定理]\label{thm:rank-nullity}
有限维时 $\dim\ker\mathbf{T}+\dim\im\mathbf{T}=\dim V$。等价地，第一同构定理给出 $V/\ker\mathbf{T}\cong\im\mathbf{T}$。
\end{theorem}
\begin{proof}
取 $\ker\mathbf{T}$ 的基 $\{\mathbf{k}_1,\dots,\mathbf{k}_p\}$，扩成 $V$ 的基 $\{\mathbf{k}_1,\dots,\mathbf{k}_p,\mathbf{e}_1,\dots,\mathbf{e}_q\}$。断言 $\{\mathbf{T}\mathbf{e}_1,\dots,\mathbf{T}\mathbf{e}_q\}$ 是 $\im\mathbf{T}$ 的基。\emph{张成}：任意 $\mathbf{T}\mathbf{v}=\mathbf{T}(\sum a_i\mathbf{k}_i+\sum b_j\mathbf{e}_j)=\sum b_j\mathbf{T}\mathbf{e}_j$（核部分被映为零）。\emph{独立}：若 $\sum b_j\mathbf{T}\mathbf{e}_j=\mathbf{0}$ 则 $\sum b_j\mathbf{e}_j\in\ker\mathbf{T}$，可由 $\mathbf{k}_i$ 表出，由基的独立性得 $b_j=0$。故 $\dim\im\mathbf{T}=q$，而 $\dim V=p+q$。
\end{proof}

\begin{example}[秩--零度即“冗余度 + 任务维数 = 关节数”]\label{ex:rank-nullity-robot}
设机械臂雅可比 $\mathbf{J}=\mathbf{J}(\boldsymbol{\theta}):\mathbb{R}^n\to\mathbb{R}^6$（$n$ 个关节速度 $\to$ 6 维末端速度旋量）。$\ker\mathbf{J}$ 是\emph{自运动}方向（关节在动而末端不动），$\im\mathbf{J}$ 是\emph{可达}末端速度方向。\cref{thm:rank-nullity} 说 $\dim\ker\mathbf{J}+\dim\im\mathbf{J}=n$，即“冗余度 $+$ 可达任务维数 $=$ 关节自由度”。奇异位形处 $\mathrm{rank}\,\mathbf{J}<6$，某些末端方向\emph{不可达}（$\dim\im\mathbf{J}$ 下降），同时冗余度 $\dim\ker\mathbf{J}$ 上升。
\end{example}

\paragraph{矩阵只是基下的影子。}
选定 $V,W$ 的基后，$\mathbf{T}$ 写成矩阵 $[\mathbf{T}]$。换基由可逆矩阵 $\mathbf{P}$ 描述，\emph{相似变换} $[\mathbf{T}]_{\mathcal{B}'}=\mathbf{P}^{-1}[\mathbf{T}]_{\mathcal{B}}\mathbf{P}$。相似矩阵代表\emph{同一}线性映射，故迹、行列式、特征值、秩都是\emph{相似不变量}——它们刻画映射本身、与坐标无关。这是“坐标无关思维”的第一课：凡是坐标系一换就变的量都不是映射的内禀属性。

\begin{practice}
\begin{enumerate}
  \item 验证 $M_{n}(\mathbb{R})$ 中对称矩阵构成子空间，求其维数（提示：$n(n+1)/2$）；反对称矩阵呢？$\so(3)$ 是哪一个的子空间？
  \item 设 $\mathbf{T}:\mathbb{R}^4\to\mathbb{R}^3$ 矩阵为 $\bigl[\begin{smallmatrix}1&2&0&1\\0&0&1&3\\1&2&1&4\end{smallmatrix}\bigr]$。求 $\ker\mathbf{T}$、$\im\mathbf{T}$ 的基，验证 \cref{thm:rank-nullity}。
  \item 解释为什么迹与行列式是相似不变量（提示：$\det(\mathbf{P}^{-1}\mathbf{A}\mathbf{P})=\det\mathbf{A}$、$\tr(\mathbf{A}\mathbf{B})=\tr(\mathbf{B}\mathbf{A})$）。
\end{enumerate}
\end{practice}

% ════════════════════════════════════════════════════════════════
\section{对偶空间：twist 与 wrench 的代数}
\label{sec:dual}

\paragraph{动机：力和速度真住在同一个空间吗？}
力 $\mathbf{F}$ 与速度 $\mathbf{v}$ 的配对 $\mathbf{F}^\top\mathbf{v}$ 给出\emph{功率}（标量）。朴素看二者都是 $\mathbb{R}^3$，但坐标一换它们\emph{行为不同}：速度按 $\mathbf{v}'=\mathbf{R}\mathbf{v}$ 变换（\emph{逆变}），力按 $\mathbf{F}'=\mathbf{R}^{-\top}\mathbf{F}$ 变换（\emph{协变}）。在正交基下 $\mathbf{R}^{-\top}=\mathbf{R}$ 看不出区别，非正交（广义坐标）时区别立现。对偶空间正是把这种区别\emph{形式化}：力住在速度空间的\emph{对偶}里。

\begin{definition}[对偶空间与对偶配对]\label{def:dual}
$V$ 上的\emph{线性泛函}是线性映射 $f:V\to\mathbb{F}$。全体线性泛函构成 $V$ 的\emph{对偶空间} $V^*:=\Hom(V,\mathbb{F})$，本身是向量空间。\emph{对偶配对} $\langle f,\mathbf{v}\rangle:=f(\mathbf{v})$。
\end{definition}

\noindent 几何上 $f\in V^*$ 是一种“测量”：它把向量读成一个数，其核 $\ker f$ 是一张\emph{超平面}（余维 1）。给定 $V$ 的基 $(\mathbf{v}_1,\dots,\mathbf{v}_n)$，\emph{对偶基} $(f^1,\dots,f^n)\subset V^*$ 由 $f^i(\mathbf{v}_j)=\delta^i_j$ 唯一确定，它“只读第 $i$ 个坐标分量”。由此 $\dim V^*=\dim V$（有限维）。

\begin{pitfall}[概念误区：$V^*\cong V$ 所以“没区别”]
$\dim V^*=\dim V$ 仅意味着二者\emph{同构}，但任何 $V\to V^*$ 的同构都要\emph{先选一组基}（或选一个内积），换基就换同构——没有“天然”的对应。$V$ 与 $V^*$ 是两个\emph{不同}空间：元素物理含义不同（速度 vs 力）、变换律不同（逆变 vs 协变）。在机器人学里，twist $\mathbf{V}\in\se(3)$ 按 $\mathrm{Ad}_{\mathbf{T}}$ 变换、wrench $\mathbf{F}\in\se(3)^*$ 按 $\mathrm{Ad}_{\mathbf{T}}^{-\top}$ 变换；混淆二者会把动力学方程写错。
\end{pitfall}

\paragraph{对偶映射（转置）。}
线性映射 $\mathbf{T}:V\to W$ 自然诱导一个\emph{反方向}的\emph{对偶映射} $\mathbf{T}^\top:W^*\to V^*$，由“把泛函沿 $\mathbf{T}$ 拉回”定义：
\begin{equation}
  (\mathbf{T}^\top g)(\mathbf{v}):=g(\mathbf{T}\mathbf{v}),\qquad g\in W^*.
  \label{eq:dual-map}
\end{equation}
在对偶基下，$\mathbf{T}^\top$ 的矩阵恰是 $\mathbf{T}$ 矩阵的\emph{转置}。注意 \cref{eq:dual-map} \textbf{不需要任何内积}——它纯由线性结构给出。

\begin{example}[$\boldsymbol{\tau}=\mathbf{J}^\top\mathbf{F}$ 是对偶映射，不是伴随]\label{ex:jacobian-transpose}
机械臂静力学的核心公式“关节力矩 $=$ 雅可比转置乘末端力” $\boldsymbol{\tau}=\mathbf{J}^\top\mathbf{F}$ 由\emph{虚功原理}得出：任意虚位移下关节做的功 $\boldsymbol{\tau}^\top\delta\boldsymbol{\theta}$ 必等于末端做的功 $\mathbf{F}^\top\delta\mathbf{x}=\mathbf{F}^\top\mathbf{J}\,\delta\boldsymbol{\theta}$，对任意 $\delta\boldsymbol{\theta}$ 成立故 $\boldsymbol{\tau}=\mathbf{J}^\top\mathbf{F}$。这里 $\mathbf{J}^\top$ 是把任务空间的 wrench（$\in W^*$）反拉回关节力矩空间（$\in V^*$）的\emph{对偶映射}，是\emph{精确等式}、\textbf{不依赖任何度量}。若错把它当“内积伴随”，就会误以为它依赖度量、在非正交关节坐标下写出错误公式。对偶映射 $\mathbf{T}^\top$ 与伴随算子 $\mathbf{A}^*$ 的区分见 \cref{sec:inner-product}。
\end{example}

\paragraph{二次对偶与“自然”同构。}
\emph{评估映射} $\mathrm{ev}:V\to V^{**}$，$\mathrm{ev}_{\mathbf{v}}(f):=f(\mathbf{v})$，把向量送成“在它处求值”的泛函。它是\emph{单射}（取一组含 $\mathbf{v}$ 的基可造出 $f$ 使 $f(\mathbf{v})\neq0$），有限维时是同构。关键差别：$V\cong V^*$ 需\emph{选基}（非自然），而 $V\cong V^{**}$ 由 $\mathrm{ev}$ 给出、\emph{对所有 $V$ 一致、不依赖任何选择}（自然）。

\begin{insight}[“自然”=帧无关，“非自然”=需要度量]
范畴论意义上 $\mathrm{ev}:\mathrm{Id}\Rightarrow(-)^{**}$ 是\emph{自然变换}，而不存在 $\mathrm{Id}\Rightarrow(-)^*$ 的自然变换（一个协变、一个反变，方差不同）。物理含义直接：twist 与它的“二次对偶”识别是\emph{帧无关}的（自然）；twist 与 wrench 的换算却\emph{需要惯性张量}（一个“内积”），是\emph{非自然}的。这解释了为何“为什么是转置”的疑惑在机器人学反复出现——一旦建立对偶直觉，$\mathbf{J}^\top$ 是 $\mathbf{J}$ 的对偶映射、coadjoint 是 adjoint 的对偶作用，大量疑惑自动消解。
\end{insight}

\paragraph{零化子与“行秩 = 列秩”。}
子空间 $W\le V$ 的\emph{零化子} $W^\circ:=\{f\in V^*:f|_W=0\}\le V^*$，满足维数公式 $\dim W+\dim W^\circ=\dim V$。它给出核--像的对偶关系
\begin{equation}
  \ker\mathbf{T}^\top=(\im\mathbf{T})^\circ,\qquad \im\mathbf{T}^\top=(\ker\mathbf{T})^\circ,
  \label{eq:annihilator-duality}
\end{equation}
由此 $\mathrm{rank}\,\mathbf{T}^\top=\dim W-\dim(\im\mathbf{T})^\circ=\mathrm{rank}\,\mathbf{T}$——这正是“\emph{行秩 $=$ 列秩}”的\emph{概念}证明（无需行化简）。在机器人学中，$W^\circ$ 翻译为“约束力 / 互易旋量（reciprocal screws）”：作用在约束切空间正交补上的 wrench 不做功。下面把这句“一句话译介”展开成一个完整的对偶结构。

\paragraph{从零化子到互易螺旋系。}
零化子 $W^\circ$ 在机器人学有个具体化身——\emph{互易螺旋系}（reciprocal screw system）。把刚体瞬时运动写成 twist $\mathbf{V}=(\mathbf{v};\boldsymbol{\omega})\in\se(3)$（本书 $\boldsymbol{\xi}=[\boldsymbol{\rho};\boldsymbol{\phi}]$ 序，平移在前），作用其上的力/力偶写成 wrench $\mathbf{F}=(\mathbf{f};\boldsymbol{\tau})\in\se(3)^*$，则功率配对 $\langle\mathbf{F},\mathbf{V}\rangle=\mathbf{v}\cdot\mathbf{f}+\boldsymbol{\omega}\cdot\boldsymbol{\tau}$（与 \cref{ex:jacobian-transpose} 的虚功同一配对）。

\begin{definition}[互易（reciprocity）与互易积]\label{def:recip-product}
wrench $\mathbf{F}$ 与 twist $\mathbf{V}$ 称\emph{互易}，若瞬时功率为零 $\langle\mathbf{F},\mathbf{V}\rangle=0$（运动不抵抗该力 / 力不阻碍该运动）。当两螺旋 $S_1,S_2$ 都用 twist 坐标 $\mathbf{V}_i=(\mathbf{v}_i;\boldsymbol{\omega}_i)$ 表示时（把其一转成 wrench 再配对），其\emph{互易积}为
\begin{equation}
  \mathbf{V}_1\odot\mathbf{V}_2=\mathbf{v}_1^\top\boldsymbol{\omega}_2+\mathbf{v}_2^\top\boldsymbol{\omega}_1=\mathbf{V}_1^\top\boldsymbol{\Pi}\,\mathbf{V}_2,\qquad \boldsymbol{\Pi}=\begin{psmallmatrix}\mathbf{0}&\mathbf{I}\\\mathbf{I}&\mathbf{0}\end{psmallmatrix},
  \label{eq:recip-product}
\end{equation}
两螺旋\emph{互易} $\iff\mathbf{V}_1\odot\mathbf{V}_2=0$。$\boldsymbol{\Pi}$（交换上下块）是\emph{对称非退化}双线性型（Klein 型），\textbf{不是}欧氏内积——互易是\emph{做功}意义的“正交”，与度量无关（恰如 $W^\circ$ 不需内积）。
\end{definition}

\begin{note}[互易积的几何形式]\label{note:recip-geom}
若把螺旋 $S_i$ 写成轴线、节距 $h_i$、幅值 $M_i$，设两轴最短距离 $d$、夹角 $\alpha$，则 $S_1\odot S_2=M_1M_2\big[(h_1+h_2)\cos\alpha-d\sin\alpha\big]$。特例：两个零节距螺旋（纯转动，$h_1=h_2=0$）互易 $\iff d\sin\alpha=0$，即两转轴\emph{共面}（相交 $d=0$ 或平行 $\sin\alpha=0$）——这给“两转动副何时互不做功”一个纯几何判据。
\end{note}

\paragraph{螺旋系与互易螺旋系。}
一组螺旋 $\{S_1,\dots,S_k\}$ 的全体线性组合（在 twist 坐标里）张成 $\se(3)$ 的子空间，称\emph{螺旋系} $\mathcal{S}$；与 $\mathcal{S}$ 中\emph{每个}螺旋都互易的全体螺旋构成\emph{互易螺旋系} $\mathcal{S}^{\perp r}$。它正是零化子在 $\se(3)$ 上的化身：

\begin{theorem}[互易螺旋系的维数：$\dim\mathcal{S}+\dim\mathcal{S}^{\perp r}=6$]\label{thm:recip-dim}
$\se(3)$ 中任一螺旋系 $\mathcal{S}$ 与其互易系 $\mathcal{S}^{\perp r}$ 满足 $\dim\mathcal{S}+\dim\mathcal{S}^{\perp r}=6$。
\end{theorem}
\begin{proof}
功率配对 $\se(3)\times\se(3)^*\to\mathbb{R}$ 非退化，故它把 $\se(3)^*$ 等同于 $\se(3)$（等价地，Klein 型 $\boldsymbol{\Pi}$ 非退化给出同构 $\se(3)\cong\se(3)^*$）。在此等同下，互易系 $\mathcal{S}^{\perp r}$ 恰是螺旋系 $\mathcal{S}$ 的零化子 $\mathcal{S}^\circ\subseteq\se(3)^*$。由零化子维数公式 $\dim\mathcal{S}+\dim\mathcal{S}^\circ=\dim\se(3)=6$（\cref{eq:annihilator-duality} 上方），即得。
\end{proof}

\begin{insight}[约束 $\leftrightarrow$ 自由度：同一对偶的两种读法]\label{ins:recip-duality}
\cref{thm:recip-dim} 把“约束多一维、自由度少一维”钉死成 $\se(3)$ 里的零化子守恒。两种读法：(i) 若 $\mathcal{S}$ 是\emph{约束 wrench 系}（如抓取中各接触点法向力、关节支反力），则 $\mathcal{S}^{\perp r}$ 是物体\emph{仍可做、不被这些力抵抗}的运动 twist——\emph{自由度}；接触每加一道独立约束，$\dim\mathcal{S}$ 加一、自由度减一，加到 $6$ 即\emph{形封闭}（form closure，物体被完全锁死）。(ii) 反过来若 $\mathcal{S}$ 是机构\emph{允许运动 twist 系}（如各关节轴），则 $\mathcal{S}^{\perp r}$ 是机构能被动承受、\emph{不做功}的约束 wrench。这正是约束动力学里 $\ker\mathbf{J}$（容许运动）与 $\operatorname{row}\mathbf{J}$（约束力方向）正交互补的螺旋几何版本——\cref{ex:jacobian-transpose} 的 $\boldsymbol{\tau}=\mathbf{J}^\top\mathbf{F}$ 对偶在此从“一个 Jacobian”升级为“一对互易子空间”。
\end{insight}

\begin{example}[转动副的自由度 twist 与其 $5$ 维互易 wrench 系]\label{ex:recip-revolute}
绕过原点、沿 $\hat{\mathbf{z}}$ 轴的转动副：自由度是纯转动 twist $\mathbf{V}=(\mathbf{v};\boldsymbol{\omega})=(\mathbf{0};\hat{\mathbf{z}})$（轴过原点故 $\mathbf{v}=-\boldsymbol{\omega}\times\mathbf{0}=\mathbf{0}$），$\dim\mathcal{S}=1$。互易 wrench $\mathbf{F}=(\mathbf{f};\boldsymbol{\tau})$ 须 $\langle\mathbf{F},\mathbf{V}\rangle=\mathbf{0}\cdot\mathbf{f}+\hat{\mathbf{z}}\cdot\boldsymbol{\tau}=\tau_z=0$。故互易系 $=\{(\mathbf{f};\boldsymbol{\tau}):\tau_z=0\}$：任意力 $\mathbf{f}$（$3$ 维）加 $z$ 分量为零的力偶 $\boldsymbol{\tau}$（$2$ 维），共 $5$ 维，验证 $\dim\mathcal{S}+\dim\mathcal{S}^{\perp r}=1+5=6$。物理上：转动副自由转、不抗任何\emph{对该轴无力矩}的 wrench（任意力、任意垂直 $\hat{\mathbf{z}}$ 的力偶皆不做功）；它唯一要抵抗的是绕 $\hat{\mathbf{z}}$ 的力矩——驱动/锁定该副所需，恰落在互易系之外。
\end{example}

\begin{practice}
\begin{enumerate}
  \item 设 $V=\mathbb{R}^3$，$f(\mathbf{x})=2x_1-x_2+3x_3$。验证 $f\in V^*$，描述 $\ker f$（一张超平面），算 $\langle f,(1,2,3)\rangle$。
  \item 平面内（$\se(2)$，$\dim=3$）取沿 $\hat{\mathbf{z}}$ 的纯转动 twist 与一条平面内纯力 wrench，用 \cref{eq:recip-product} 判其是否互易；并验证 $\se(2)$ 中螺旋系与互易系的维数和为 $3$（即 \cref{thm:recip-dim} 的平面版）。
  \item 设 $\mathbf{T}:\mathbb{R}^3\to\mathbb{R}^2$，$\mathbf{T}(x_1,x_2,x_3)=(x_1+x_2,x_2+x_3)$。写出 $\mathbf{T}$ 在标准基的矩阵 $\mathbf{A}$，验证 $\mathbf{T}^\top$ 在对偶基的矩阵是 $\mathbf{A}^\top$，并说明 $\mathbf{T}^\top$ 方向为 $(\mathbb{R}^2)^*\to(\mathbb{R}^3)^*$。
  \item 用 \cref{eq:annihilator-duality} 证明 $\mathrm{rank}\,\mathbf{T}=\mathrm{rank}\,\mathbf{T}^\top$。
\end{enumerate}
\end{practice}

% ════════════════════════════════════════════════════════════════
\section{内积空间、正交投影与伴随算子}
\label{sec:inner-product}

\paragraph{动机：给裸空间装上“长度”和“角度”。}
\cref{sec:vector-space,sec:dual} 的向量空间没有“长度”“角度”“垂直”——而机器人的几乎每个问题都需要它们：状态估计问“估计离真值多远”（距离）、控制问“姿态偏离多少”（角度/模长）、SLAM 最小化残差（范数）、点云配准最大化对齐（内积）。这些是\emph{额外}结构，必须主动装上去。装上的这台“度量仪器”叫\emph{内积}。

\begin{definition}[内积空间]\label{def:inner-product}
设 $V$ 是 $\mathbb{R}$ 或 $\mathbb{C}$ 上的向量空间。\emph{内积} $\langle\cdot,\cdot\rangle:V\times V\to\mathbb{F}$ 满足：(i) 第一变量线性 $\langle a\mathbf{u}+\mathbf{w},\mathbf{v}\rangle=a\langle\mathbf{u},\mathbf{v}\rangle+\langle\mathbf{w},\mathbf{v}\rangle$；(ii) 共轭对称 $\langle\mathbf{u},\mathbf{v}\rangle=\overline{\langle\mathbf{v},\mathbf{u}\rangle}$（实情形即对称）；(iii) 正定 $\langle\mathbf{v},\mathbf{v}\rangle\ge0$ 且 $=0\iff\mathbf{v}=\mathbf{0}$。诱导范数 $\|\mathbf{v}\|:=\sqrt{\langle\mathbf{v},\mathbf{v}\rangle}$。
\end{definition}

\noindent 复情形的“共轭”不是装饰，而是\emph{正定性的守护者}：若硬用双线性对称，则 $\langle i\mathbf{v},i\mathbf{v}\rangle=i^2\langle\mathbf{v},\mathbf{v}\rangle=-\langle\mathbf{v},\mathbf{v}\rangle<0$，长度平方变负；共轭使 $\langle i\mathbf{v},i\mathbf{v}\rangle=i\bar{i}\langle\mathbf{v},\mathbf{v}\rangle=\langle\mathbf{v},\mathbf{v}\rangle\ge0$。标准例子：$\mathbb{R}^n$ 点积 $\mathbf{x}^\top\mathbf{y}$；矩阵 Frobenius 内积 $\langle\mathbf{A},\mathbf{B}\rangle=\tr(\mathbf{A}\mathbf{B}^\top)$；\emph{加权/能量内积} $\langle\mathbf{x},\mathbf{y}\rangle_{\mathbf{M}}=\mathbf{x}^\top\mathbf{M}\mathbf{y}$（$\mathbf{M}$ 对称正定）。

\begin{insight}[同一空间可装不同内积——“用哪个度量”是物理问题]
内积是\emph{外加、可选}的：$\mathbb{R}^n$ 上标准点积、加权 $\mathbf{x}^\top\mathbf{M}\mathbf{y}$、能量内积各给一套长度与角度。机器人学里“用哪个内积”往往是\emph{物理}而非数学问题：机械臂动能 $T=\tfrac12\dot{\mathbf{q}}^\top\mathbf{M}(\mathbf{q})\dot{\mathbf{q}}$ 中的质量矩阵 $\mathbf{M}(\mathbf{q})$ 就在关节速度空间定义了“\emph{动能内积}”，动态一致伪逆用的正是它。把关节空间当欧氏空间做插值（忽略不同关节惯量差异）就是误用了度量。
\end{insight}

\paragraph{Cauchy--Schwarz、正交与勾股。}
内积满足 \emph{Cauchy--Schwarz} 不等式 $|\langle\mathbf{u},\mathbf{v}\rangle|\le\|\mathbf{u}\|\,\|\mathbf{v}\|$（等号当且仅当线性相关），它保证 $\cos\theta=\langle\mathbf{u},\mathbf{v}\rangle/(\|\mathbf{u}\|\|\mathbf{v}\|)\in[-1,1]$，使“夹角”有定义，并由此推出三角不等式。$\langle\mathbf{u},\mathbf{v}\rangle=0$ 称\emph{正交}，此时\emph{勾股}成立 $\|\mathbf{u}+\mathbf{v}\|^2=\|\mathbf{u}\|^2+\|\mathbf{v}\|^2$。

\begin{insight}[Cauchy--Schwarz 的变分证明就是一维最小二乘]
取 $\mathbf{v}\neq\mathbf{0}$，对实标量 $t$ 由 $\|\mathbf{u}-t\mathbf{v}\|^2=\|\mathbf{u}\|^2-2t\langle\mathbf{u},\mathbf{v}\rangle+t^2\|\mathbf{v}\|^2\ge0$，这是开口向上的二次函数，最小值在 $t^*=\langle\mathbf{u},\mathbf{v}\rangle/\|\mathbf{v}\|^2$ 处，代入得 $\|\mathbf{u}\|^2-\langle\mathbf{u},\mathbf{v}\rangle^2/\|\mathbf{v}\|^2\ge0$ 即不等式。这里“找 $t$ 使 $\|\mathbf{u}-t\mathbf{v}\|$ 最小”正是\emph{一维最小二乘}：$\mathbf{u}$ 是观测、$\mathbf{v}$ 是模型基、$t^*$ 是估计。把 $\mathbf{v}$ 换成一组基（矩阵 $\mathbf{A}$）就得一般最小二乘——所以这个证明是机器人状态估计的数学胚胎。
\end{insight}

\paragraph{Gram--Schmidt、正交补与正交投影。}
任一线性无关组可经 \emph{Gram--Schmidt} 正交化（逐个减去在已得方向上的投影），给出标准正交基，并等价于矩阵的 \emph{QR 分解} $\mathbf{A}=\mathbf{Q}\mathbf{R}$（$\mathbf{Q}$ 列正交、$\mathbf{R}$ 上三角）。子空间 $W$ 的\emph{正交补} $W^\perp=\{\mathbf{v}:\langle\mathbf{v},\mathbf{w}\rangle=0,\forall\mathbf{w}\in W\}$ 给出\emph{正交分解} $V=W\oplus W^\perp$，这是 \cref{sec:vector-space} 中代数补的\emph{特例}——额外要求两补相互正交，从而\emph{唯一}。

\begin{theorem}[正交投影与法方程]\label{thm:projection}
设 $W=\im\mathbf{A}\le V$，$\mathbf{A}$ 列满秩。$\mathbf{v}$ 到 $W$ 的\emph{最佳逼近}（最小化 $\|\mathbf{v}-\mathbf{w}\|$，$\mathbf{w}\in W$）是正交投影 $\proj_W\mathbf{v}=\mathbf{A}(\mathbf{A}^\top\mathbf{A})^{-1}\mathbf{A}^\top\mathbf{v}$，其特征是\emph{残差正交于 $W$}：$\langle\mathbf{v}-\proj_W\mathbf{v},\mathbf{w}\rangle=0\ \forall\mathbf{w}\in W$。等价地，$\min_{\mathbf{x}}\|\mathbf{A}\mathbf{x}-\mathbf{v}\|^2$ 的解满足\emph{法方程}
\begin{equation}
  \mathbf{A}^\top\mathbf{A}\,\mathbf{x}=\mathbf{A}^\top\mathbf{v}.
  \label{eq:normal-equation}
\end{equation}
\end{theorem}
\begin{proof}
设 $\mathbf{w}^\star=\proj_W\mathbf{v}$ 由残差正交性定义，$\mathbf{w}\in W$ 任意。$\|\mathbf{v}-\mathbf{w}\|^2=\|(\mathbf{v}-\mathbf{w}^\star)+(\mathbf{w}^\star-\mathbf{w})\|^2=\|\mathbf{v}-\mathbf{w}^\star\|^2+\|\mathbf{w}^\star-\mathbf{w}\|^2$（勾股，因 $\mathbf{v}-\mathbf{w}^\star\perp W\ni\mathbf{w}^\star-\mathbf{w}$）$\ge\|\mathbf{v}-\mathbf{w}^\star\|^2$，故 $\mathbf{w}^\star$ 最优。残差正交性 $\mathbf{A}^\top(\mathbf{v}-\mathbf{A}\mathbf{x})=\mathbf{0}$ 即 \cref{eq:normal-equation}。
\end{proof}

\begin{insight}[法方程 = 把观测正交投影到模型列空间]
\cref{eq:normal-equation} 不是凭空的代数操作，而是\emph{几何}的：最小二乘解 $\mathbf{x}$ 使 $\mathbf{A}\mathbf{x}$ 是 $\mathbf{v}$ 在 $\im\mathbf{A}$ 上的正交投影，残差 $\perp$ 列空间。SLAM 后端、ICP、束调整里反复出现的 $\mathbf{H}^\top\boldsymbol{\Sigma}^{-1}\mathbf{H}\,\delta\mathbf{x}=\mathbf{H}^\top\boldsymbol{\Sigma}^{-1}\mathbf{r}$ 就是\emph{加权}法方程（用信息矩阵 $\boldsymbol{\Sigma}^{-1}$ 定义的内积）。卡尔曼增益 $\mathbf{K}=\mathbf{P}\mathbf{H}^\top(\mathbf{H}\mathbf{P}\mathbf{H}^\top+\mathbf{R})^{-1}$ 是它的贝叶斯版本，条件期望 = 正交投影、新息序列 = 测量的 Gram--Schmidt 正交化、协方差更新 = 勾股定理。\cref{eq:normal-equation} 的左端 $\mathbf{A}^\top\mathbf{A}$ 对称正定，故可 Cholesky 分解 $\mathbf{A}^\top\mathbf{A}=\mathbf{L}\mathbf{L}^\top$（P1 优化的数值核心，见 \cref{sec:spectral} 的对称正定平方根）。
\end{insight}

\paragraph{Riesz 表示与伴随算子。}
有限维内积空间里，每个泛函都由一个向量“代表”：

\begin{theorem}[Riesz 表示定理]\label{thm:riesz}
有限维内积空间 $V$ 中，映射 $\mathbf{v}\mapsto\langle\cdot,\mathbf{v}\rangle$ 是 $V\to V^*$ 的（共轭线性）同构。即每个 $\varphi\in V^*$ 唯一地写成 $\varphi(\cdot)=\langle\cdot,\mathbf{u}\rangle$。
\end{theorem}

\noindent Riesz 把“对偶”与“内积”两条平行结构焊接起来，由此可定义\emph{伴随算子}：对 $\mathbf{A}:V\to W$，存在唯一 $\mathbf{A}^*:W\to V$ 使
\begin{equation}
  \langle\mathbf{A}\mathbf{x},\mathbf{y}\rangle_W=\langle\mathbf{x},\mathbf{A}^*\mathbf{y}\rangle_V,\qquad\forall\mathbf{x},\mathbf{y}.
  \label{eq:adjoint}
\end{equation}
标准正交基下 $\mathbf{A}^*=\bar{\mathbf{A}}^\top$（实情形 $=\mathbf{A}^\top$）；\emph{一般内积}下需经 Gram 矩阵做度规修正。

\begin{pitfall}[术语警示：伴随 $\mathbf{A}^*$、对偶 $\mathbf{T}^\top$、李群伴随 $\mathrm{Ad}$ 是三件事]
机器人学里“adjoint / 转置 / 伴随”极易混。\textbf{(1) 对偶映射} $\mathbf{T}^\top:W^*\to V^*$（\cref{eq:dual-map}）纯代数、\emph{不需内积}，矩阵为转置 $\mathbf{A}^\top$。\textbf{(2) 内积伴随} $\mathbf{A}^*:W\to V$（\cref{eq:adjoint}）\emph{需内积}，矩阵为共轭转置 $\bar{\mathbf{A}}^\top$。二者仅在“实数域 $+$ 标准正交基”这一三重巧合下矩阵相同。\textbf{(3) 李群伴随} $\mathrm{Ad}_{\mathbf{g}}:\mathfrak{g}\to\mathfrak{g}$，$\boldsymbol{\xi}\mapsto\mathbf{g}\boldsymbol{\xi}\mathbf{g}^{-1}$（\cref{ch:lie}），与前两者同名而\emph{异义}。三者在 wrench 变换 $\mathbf{F}_b=[\mathrm{Ad}_{\mathbf{T}}]^\top\mathbf{F}_a$ 中\emph{合一}：$[\mathrm{Ad}_{\mathbf{T}}]^\top$ 既是 coadjoint（作用在对偶 $\se(3)^*$），又是 $[\mathrm{Ad}_{\mathbf{T}}]$ 在标准内积下的伴随。协方差传播 $\boldsymbol{\Sigma}'=[\mathrm{Ad}_{\mathbf{X}}]\boldsymbol{\Sigma}[\mathrm{Ad}_{\mathbf{X}}]^\top$ 同时用到两种伴随——这是全书“结构统一性”最精彩的体现。
\end{pitfall}

\paragraph{算子分类：自伴、正规、酉。}
伴随把算子分成三类（谱定理的入场券）：\emph{自伴} $\mathbf{A}=\mathbf{A}^*$（实即对称 $\mathbf{A}=\mathbf{A}^\top$）；\emph{正规} $\mathbf{A}\mathbf{A}^*=\mathbf{A}^*\mathbf{A}$；\emph{酉/正交} $\mathbf{A}^*\mathbf{A}=\mathbf{I}$（保内积、保长度，逆 $=$ 共轭转置）。包含关系：自伴 $\subset$ 正规、酉 $\subset$ 正规，但自伴与酉互不包含。

\begin{example}[$\mathbf{R}^{-1}=\mathbf{R}^\top$ 不是巧合，是“正交=保内积”的必然]\label{ex:orthogonal}
旋转 $\mathbf{R}\in\SO(3)$ 保长度（保范），即对任意 $\mathbf{x},\mathbf{y}$ 有 $\langle\mathbf{R}\mathbf{x},\mathbf{R}\mathbf{y}\rangle=\langle\mathbf{x},\mathbf{y}\rangle$。由 \cref{eq:adjoint}，$\langle\mathbf{x},\mathbf{R}^\top\mathbf{R}\mathbf{y}\rangle=\langle\mathbf{x},\mathbf{y}\rangle$ 对所有 $\mathbf{x}$ 成立，故 $\mathbf{R}^\top\mathbf{R}=\mathbf{I}$ 即 $\mathbf{R}^{-1}=\mathbf{R}^\top$。所以“旋转矩阵的逆恰是转置”是\emph{保内积}的代数后果，而非记忆性巧合（对照 \cref{ch:rigid_body}）。$\SE(3)$ 不是正交群（平移项额外处理），其逆不是简单转置。
\end{example}

\begin{practice}
\begin{enumerate}
  \item 在 $\mathcal{P}_2(\mathbb{R})$ 上验证 $\langle p,q\rangle=\int_{-1}^1 p(t)q(t)\,dt$ 是内积，算 $\langle1,t\rangle,\langle1,t^2\rangle,\langle t,t^2\rangle$，指出哪些对正交（Legendre 多项式的起点）。
  \item 推导加权伪逆 $\mathbf{J}^+_{\mathbf{W}}=\mathbf{W}^{-1}\mathbf{J}^\top(\mathbf{J}\mathbf{W}^{-1}\mathbf{J}^\top)^{-1}$ 是 $\min\dot{\boldsymbol{\theta}}^\top\mathbf{W}\dot{\boldsymbol{\theta}}$ s.t. $\mathbf{J}\dot{\boldsymbol{\theta}}=\mathbf{V}_d$ 的解（提示：Lagrange/KKT，用 $\mathbf{W}$-内积）。
  \item 证明自伴 $\subset$ 正规、酉 $\subset$ 正规，并举一个正规但既非自伴也非酉的实例（提示：$\bigl[\begin{smallmatrix}1&-1\\1&1\end{smallmatrix}\bigr]$，验证 $\mathbf{A}^\top\mathbf{A}=\mathbf{A}\mathbf{A}^\top=2\mathbf{I}$ 故正规，而 $\mathbf{A}\neq\mathbf{A}^\top$、$\mathbf{A}^\top\mathbf{A}\neq\mathbf{I}$ 故二者皆非）。\footnote{此例须取 $\mathbf{A}^\top\mathbf{A}\neq\mathbf{I}$ 者方可同时排除酉性：纯旋转 $\bigl[\begin{smallmatrix}0&-1\\1&0\end{smallmatrix}\bigr]$ 虽正规且非对称，但其满足 $\mathbf{A}^\top\mathbf{A}=\mathbf{I}$ 恰是正交（酉）矩阵，不能用作"非酉"的反例。缩放后的旋转 $\bigl[\begin{smallmatrix}1&-1\\1&1\end{smallmatrix}\bigr]=\sqrt2\,\bigl[\begin{smallmatrix}\cos45^\circ&-\sin45^\circ\\\sin45^\circ&\cos45^\circ\end{smallmatrix}\bigr]$ 把模长改成 $\sqrt2$，破坏酉性而保正规。}
\end{enumerate}
\end{practice}

% ════════════════════════════════════════════════════════════════
\section{谱定理与正定性}
\label{sec:spectral}

\paragraph{动机：把耦合的变换拆成互不干扰的伸缩。}
惯性张量 $\mathbf{I}$ 是 $3\times3$ 对称矩阵：一般绕任意轴施力矩，角加速度方向\emph{不}与力矩平行（$\boldsymbol{\tau}=\mathbf{I}\dot{\boldsymbol{\omega}}$，$\mathbf{I}$ 会“扭转”方向），动力学方程耦合。但若找到 $\mathbf{I}$ 的特征向量——\emph{主惯性轴}——绕主轴转动时力矩与角加速度平行，三轴方程彻底解耦。每本刚体力学先求主轴，正是为把耦合问题拆成独立的简单问题。我们真正想要的不是随便对角化，而是\emph{正交}对角化（特征向量互相垂直），它对应物理上垂直的转动轴、且数值稳定（基变换是正交矩阵，逆即转置）。

\begin{theorem}[实对称谱定理]\label{thm:spectral-real}
设 $\mathbf{A}\in\mathbb{R}^{n\times n}$ 对称（$\mathbf{A}=\mathbf{A}^\top$）。则 (i) 特征值全为实数；(ii) 不同特征值的特征向量\emph{自动}正交；(iii) 存在正交矩阵 $\mathbf{Q}$（$\mathbf{Q}^\top\mathbf{Q}=\mathbf{I}$）使
\begin{equation}
  \mathbf{A}=\mathbf{Q}\boldsymbol{\Lambda}\mathbf{Q}^\top,\qquad\boldsymbol{\Lambda}=\diag(\lambda_1,\dots,\lambda_n),\ \lambda_i\in\mathbb{R}.
  \label{eq:spectral}
\end{equation}
\end{theorem}
\begin{proof}
\emph{(i) 实特征值}：设 $\mathbf{A}\mathbf{v}=\lambda\mathbf{v}$（复化），$\mathbf{v}\neq\mathbf{0}$。$\langle\mathbf{A}\mathbf{v},\mathbf{v}\rangle=\lambda\|\mathbf{v}\|^2$；又 $\langle\mathbf{A}\mathbf{v},\mathbf{v}\rangle=\langle\mathbf{v},\mathbf{A}\mathbf{v}\rangle=\overline{\langle\mathbf{A}\mathbf{v},\mathbf{v}\rangle}$（自伴 $+$ 共轭对称），故 $\langle\mathbf{A}\mathbf{v},\mathbf{v}\rangle$ 实，$\lambda=\langle\mathbf{A}\mathbf{v},\mathbf{v}\rangle/\|\mathbf{v}\|^2$ 实。\emph{(ii) 正交}：$\mathbf{A}\mathbf{v}_1=\lambda_1\mathbf{v}_1$、$\mathbf{A}\mathbf{v}_2=\lambda_2\mathbf{v}_2$、$\lambda_1\neq\lambda_2$。$\lambda_1\langle\mathbf{v}_1,\mathbf{v}_2\rangle=\langle\mathbf{A}\mathbf{v}_1,\mathbf{v}_2\rangle=\langle\mathbf{v}_1,\mathbf{A}\mathbf{v}_2\rangle=\lambda_2\langle\mathbf{v}_1,\mathbf{v}_2\rangle$，故 $(\lambda_1-\lambda_2)\langle\mathbf{v}_1,\mathbf{v}_2\rangle=0\Rightarrow\langle\mathbf{v}_1,\mathbf{v}_2\rangle=0$。\emph{(iii) 完整正交基}：对 $\dim V$ 归纳。取单位特征向量 $\mathbf{q}_1$，其正交补 $U=\mathbf{q}_1^\perp$ 是 $\mathbf{A}$-不变（对 $\mathbf{w}\in U$，$\langle\mathbf{A}\mathbf{w},\mathbf{q}_1\rangle=\langle\mathbf{w},\mathbf{A}\mathbf{q}_1\rangle=\lambda_1\langle\mathbf{w},\mathbf{q}_1\rangle=0$，用了自伴），$\mathbf{A}|_U$ 仍自伴、$\dim U=n-1$，归纳得 $U$ 的正交特征基，拼上 $\mathbf{q}_1$ 即得。
\end{proof}

\noindent 三步的灵魂是同一招：\emph{用 $\mathbf{A}=\mathbf{A}^*$ 把伴随挪来挪去}。这个“正交补归纳”模板在 \cref{sec:svd} 证 SVD 时原样再现。几何上 \cref{eq:spectral} 读作“\emph{转到主轴 $\to$ 各轴独立缩放 $\to$ 转回来}”：交叉项消失、椭球摆正。非对称矩阵一般做不到正交对角化；推广到\emph{复正规}算子 $\mathbf{A}\mathbf{A}^*=\mathbf{A}^*\mathbf{A}$，则有酉对角化 $\mathbf{A}=\mathbf{U}\boldsymbol{\Lambda}\mathbf{U}^*$；对\emph{任意}复方阵，至少有 \emph{Schur 三角化} $\mathbf{A}=\mathbf{U}\mathbf{T}\mathbf{U}^*$（$\mathbf{T}$ 上三角，正规情形 $\mathbf{T}$ 退化为对角），它是非正规情形的核心工具，也是 \cref{thm:cayley-hamilton} 的证明基础。

\begin{theorem}[正定性的特征值判据]\label{thm:positive-definite}
对称 $\mathbf{A}$ 正定（$\mathbf{x}^\top\mathbf{A}\mathbf{x}>0,\forall\mathbf{x}\neq\mathbf{0}$）$\iff$ 所有 $\lambda_i>0$；半正定 $\iff$ 所有 $\lambda_i\ge0$。
\end{theorem}
\begin{proof}
由 \cref{eq:spectral} 换元 $\mathbf{y}=\mathbf{Q}^\top\mathbf{x}$（正交、$\|\mathbf{y}\|=\|\mathbf{x}\|$）：$\mathbf{x}^\top\mathbf{A}\mathbf{x}=\mathbf{y}^\top\boldsymbol{\Lambda}\mathbf{y}=\sum_i\lambda_iy_i^2$。在主轴坐标系下二次型成为无交叉项的平方和，故恒正 $\iff$ 全部 $\lambda_i>0$（任一 $\lambda_k\le0$ 取 $\mathbf{y}=\mathbf{e}_k$ 即反例）。
\end{proof}

\begin{insight}[正定 = “这个二次型的内在形状是个碗”，与坐标无关]
正定 $\iff$ 特征值全正，深层含义是：二次型的“正负性”是\emph{旋转不变量}，只取决于谱、与坐标系无关。无论怎么旋转坐标，碗就是碗、鞍就是鞍。这解释了为何正定性在机器人学如此根本：协方差矩阵\emph{半正定}（方差不能为负），优化目标在极小点的海森\emph{半正定}，Lyapunov 函数 $V(\mathbf{x})=\mathbf{x}^\top\mathbf{P}\mathbf{x}$ \emph{正定}才能证稳定——本质都在断言“这个二次型的内在形状是个碗”。这些物理量\emph{天生对称}（惯性张量来自二次积分、协方差来自期望、海森因 $\partial^2 f/\partial x_i\partial x_j$ 对称），所以谱定理才能直接套用。
\end{insight}

\begin{corollary}[对称半正定的唯一平方根 / Cholesky]\label{cor:sqrt}
对称半正定 $\mathbf{A}=\mathbf{Q}\boldsymbol{\Lambda}\mathbf{Q}^\top$（$\lambda_i\ge0$）有唯一对称半正定平方根 $\mathbf{A}^{1/2}=\mathbf{Q}\boldsymbol{\Lambda}^{1/2}\mathbf{Q}^\top$（对特征值开方，$\boldsymbol{\Lambda}^{1/2}=\diag(\sqrt{\lambda_i})$）。对称\emph{正定}矩阵还有唯一的 Cholesky 分解 $\mathbf{A}=\mathbf{L}\mathbf{L}^\top$（$\mathbf{L}$ 下三角、对角正）。
\end{corollary}

\begin{pitfall}[陷阱：矩阵平方根必须经谱分解，绝非逐元素开方]
$\mathbf{A}^{1/2}$ 不是把 $\mathbf{A}$ 的每个元素开方！必须先谱分解、对\emph{特征值}开方再变回。逐元素开方既不对称半正定、平方也不等于 $\mathbf{A}$。这个平方根是 \cref{sec:svd} 极分解、状态估计“白化”（whitening，把协方差 $\boldsymbol{\Sigma}$ 化为单位阵 $\boldsymbol{\Sigma}^{-1/2}$）的关键。Cholesky $\mathbf{L}\mathbf{L}^\top$ 比谱分解更省（$O(n^3/3)$），是 P1 优化求解法方程 \cref{eq:normal-equation}、采样多元高斯 $\mathcal{N}(\boldsymbol{\mu},\boldsymbol{\Sigma})$（取 $\mathbf{x}=\boldsymbol{\mu}+\mathbf{L}\mathbf{z}$，$\mathbf{z}\sim\mathcal{N}(\mathbf{0},\mathbf{I})$）的标准工具。
\end{pitfall}

\begin{example}[三个谱分解化身：惯性主轴、不确定性椭球、海森]\label{ex:spectral-robot}
\begin{itemize}
  \item \textbf{惯性张量主轴}：$\mathbf{I}=\mathbf{R}\diag(I_1,I_2,I_3)\mathbf{R}^\top$，$\mathbf{R}\in\SO(3)$ 三列是主惯性轴，$I_k$ 是主惯性矩——欧拉方程解耦的前提。
  \item \textbf{协方差不确定性椭球}：$\boldsymbol{\Sigma}=\mathbf{Q}\boldsymbol{\Lambda}\mathbf{Q}^\top$，$\mathbf{Q}$ 列是不确定性主方向，$\sqrt{\lambda_i}$ 是该方向标准差，$\det\boldsymbol{\Sigma}=\prod\lambda_i$ 衡量总不确定性体积（主动探索的 D-最优性指标）。
  \item \textbf{海森与优化局部形状}：$\mathbf{H}=\mathbf{Q}\boldsymbol{\Lambda}\mathbf{Q}^\top$，特征值告诉哪些方向“陡”（大 $\lambda$，收敛快）、哪些“平”（小 $\lambda$，病态），条件数 $\lambda_{\max}/\lambda_{\min}$ 决定梯度下降收敛速度。
\end{itemize}
\end{example}

\begin{practice}
\begin{enumerate}
  \item 设对称 $\mathbf{A}$ 特征值 $\lambda_1,\dots,\lambda_n$。证 $\tr\mathbf{A}=\sum\lambda_i$、$\det\mathbf{A}=\prod\lambda_i$、$\mathbf{A}$ 可逆 $\iff$ 全 $\lambda_i\neq0$。哪些对非对称矩阵仍成立？
  \item 谱定理第二步只证“异特征值正交”。说明同一特征值的特征空间内部可用 Gram--Schmidt 正交化，从而凑齐完整正交特征基（补全归纳证明中重特征值的细节）。
  \item 用 \cref{thm:positive-definite} 证明：对称正定矩阵之和仍正定；$\mathbf{B}^\top\mathbf{B}$ 总半正定，列满秩时正定。
\end{enumerate}
\end{practice}

% ════════════════════════════════════════════════════════════════
\section{奇异值分解、极分解与伪逆}
\label{sec:svd}

\paragraph{动机：当矩阵不是方阵，“特征值”失效。}
机械臂雅可比 $\mathbf{J}:\mathbb{R}^7\to\mathbb{R}^6$ 把 7 维映成 6 维，“$\mathbf{J}\dot{\mathbf{q}}=\lambda\dot{\mathbf{q}}$”毫无意义（两边维数不同），\emph{矩形矩阵没有特征值}。但我们仍想知道：哪个方向放大最厉害、哪个被压成零、输入主方向如何对应输出主方向？SVD 放弃“输入方向 $=$ 输出方向”的奢望，改为找输入空间的正交基 $\{\mathbf{v}_i\}$ 与输出空间的正交基 $\{\mathbf{u}_i\}$ 及非负数 $\sigma_i$ 使 $\mathbf{A}\mathbf{v}_i=\sigma_i\mathbf{u}_i$——“第 $i$ 个输入主方向被干净映到第 $i$ 个输出主方向、只做一次纯伸缩”。

\begin{theorem}[奇异值分解（SVD）]\label{thm:svd}
任意 $\mathbf{A}\in\mathbb{R}^{m\times n}$（秩 $r$）可分解为
\begin{equation}
  \mathbf{A}=\mathbf{U}\boldsymbol{\Sigma}\mathbf{V}^\top,
  \label{eq:svd}
\end{equation}
其中 $\mathbf{U}\in\mathbb{R}^{m\times m}$、$\mathbf{V}\in\mathbb{R}^{n\times n}$ 正交，$\boldsymbol{\Sigma}\in\mathbb{R}^{m\times n}$ “对角” $\Sigma_{ii}=\sigma_i$，$\sigma_1\ge\cdots\ge\sigma_r>0=\sigma_{r+1}=\cdots$。等价逐方向：$\mathbf{A}\mathbf{v}_i=\sigma_i\mathbf{u}_i\,(i\le r)$、$\mathbf{A}\mathbf{v}_i=\mathbf{0}\,(i>r)$。
\end{theorem}
\begin{proof}
命门：$\mathbf{A}^\top\mathbf{A}$（$n\times n$）\emph{永远}对称半正定——对称显然，半正定因 $\mathbf{x}^\top\mathbf{A}^\top\mathbf{A}\mathbf{x}=\|\mathbf{A}\mathbf{x}\|^2\ge0$。\textbf{Step 1}：对 $\mathbf{A}^\top\mathbf{A}$ 用谱定理（\cref{thm:spectral-real}），得正交特征基 $\mathbf{v}_1,\dots,\mathbf{v}_n$ 与非负特征值，记为 $\sigma_i^2$（降序），定义\emph{奇异值} $\sigma_i=\sqrt{\sigma_i^2}$、\emph{右奇异向量} $\mathbf{v}_i$。由 $\|\mathbf{A}\mathbf{v}_i\|^2=\mathbf{v}_i^\top\mathbf{A}^\top\mathbf{A}\mathbf{v}_i=\sigma_i^2$ 知 $\sigma_i=0\iff\mathbf{A}\mathbf{v}_i=\mathbf{0}\iff\mathbf{v}_i\in\ker\mathbf{A}$，故正奇异值个数 $r=\mathrm{rank}\,\mathbf{A}$（秩--零度，\cref{thm:rank-nullity}）。\textbf{Step 2}：对 $i\le r$ 定义\emph{左奇异向量} $\mathbf{u}_i:=\mathbf{A}\mathbf{v}_i/\sigma_i$。其正交性\emph{自动}得到：$\langle\mathbf{u}_i,\mathbf{u}_j\rangle=\frac{1}{\sigma_i\sigma_j}\langle\mathbf{A}\mathbf{v}_i,\mathbf{A}\mathbf{v}_j\rangle=\frac{1}{\sigma_i\sigma_j}\langle\mathbf{v}_i,\mathbf{A}^\top\mathbf{A}\mathbf{v}_j\rangle=\frac{\sigma_j}{\sigma_i}\delta_{ij}=\delta_{ij}$。\textbf{Step 3}：用 Gram--Schmidt 把 $\{\mathbf{u}_1,\dots,\mathbf{u}_r\}$ 扩成 $\mathbb{R}^m$ 的正交基（补进的张成 $\im\mathbf{A}$ 的正交补 $=\ker\mathbf{A}^\top$）。\textbf{Step 4}：摆成 $\mathbf{U},\mathbf{V}$，逐列验证 $\mathbf{A}\mathbf{V}=\mathbf{U}\boldsymbol{\Sigma}$（第 $i$ 列两边都是 $\sigma_i\mathbf{u}_i$ 或 $\mathbf{0}$），右乘 $\mathbf{V}^\top$ 得 \cref{eq:svd}。
\end{proof}

\noindent 灵魂一句话：\emph{一般 $\mathbf{A}$ 不好对角化，但 $\mathbf{A}^\top\mathbf{A}$ 一定对称半正定、可套谱定理}——这是“先讲谱定理再讲 SVD”的因果链。

\begin{figure}[ht]
\centering
\begin{tikzpicture}[scale=1.0, >=Stealth]
  % unit circle (input)
  \begin{scope}[shift={(-5,0)}]
    \draw[thick, main] (0,0) circle (0.9);
    \draw[->, gray] (-1.2,0)--(1.2,0); \draw[->, gray] (0,-1.2)--(0,1.2);
    \draw[->, second, thick] (0,0)--(0.9,0) node[below right,font=\footnotesize] {$\mathbf{v}_1$};
    \draw[->, second, thick] (0,0)--(0,0.9) node[above left,font=\footnotesize] {$\mathbf{v}_2$};
    \node[font=\small] at (0,-1.6) {输入单位球};
  \end{scope}
  % arrow + label
  \draw[->, very thick, third] (-3.7,0)--(-2.1,0) node[midway,above,font=\footnotesize] {$\mathbf{A}=\mathbf{U}\boldsymbol{\Sigma}\mathbf{V}^\top$};
  % ellipse (output)
  \begin{scope}[shift={(0,0)}, rotate=28]
    \draw[thick, main] (0,0) ellipse (1.5 and 0.6);
    \draw[->, second, thick] (0,0)--(1.5,0) node[above,font=\footnotesize] {$\sigma_1\mathbf{u}_1$};
    \draw[->, second, thick] (0,0)--(0,0.6) node[left,font=\footnotesize] {$\sigma_2\mathbf{u}_2$};
  \end{scope}
  \node[font=\small] at (0,-1.6) {输出椭球（半轴 $\sigma_i$，主轴 $\mathbf{u}_i$）};
\end{tikzpicture}
\caption{SVD 的几何：任意线性映射把输入单位球映成椭球。$\mathbf{V}^\top$ 先把球转一下（不变形）$\to$ $\boldsymbol{\Sigma}$ 沿坐标轴各自伸缩 $\sigma_i$（球变轴对齐椭球）$\to$ $\mathbf{U}$ 把椭球转到最终方位（不变形）。半轴长 $=$ 奇异值，主轴方向 $=$ 左奇异向量。}
\label{fig:svd-ellipse}
\end{figure}

\paragraph{几何：旋转--伸缩--旋转，四个基本子空间。}
\cref{eq:svd} 读作三段（从右往左作用）：$\mathbf{V}^\top$ 输入侧换正交坐标（不变形）$\to$ $\boldsymbol{\Sigma}$ 沿坐标轴各自伸缩 $\sigma_i$（全部形变在此）$\to$ $\mathbf{U}$ 输出侧换正交坐标（不变形）。故单位球必被映成椭球（\cref{fig:svd-ellipse}）。SVD 还\emph{一次性}给出矩阵的全部子空间结构：

\begin{table}[H]\centering\small
\caption{SVD 给出四个基本子空间的正交基（$\mathbf{A}=\mathbf{U}\boldsymbol{\Sigma}\mathbf{V}^\top$，秩 $r$）}
\begin{tabular}{llll}
\toprule
子空间 & 记号 & 维数 & SVD 正交基 \\
\midrule
行空间 & $\im\mathbf{A}^\top\subseteq\mathbb{R}^n$ & $r$ & $\{\mathbf{v}_1,\dots,\mathbf{v}_r\}$（正奇异值） \\
零空间 & $\ker\mathbf{A}\subseteq\mathbb{R}^n$ & $n-r$ & $\{\mathbf{v}_{r+1},\dots,\mathbf{v}_n\}$（零奇异值） \\
列空间 & $\im\mathbf{A}\subseteq\mathbb{R}^m$ & $r$ & $\{\mathbf{u}_1,\dots,\mathbf{u}_r\}$ \\
左零空间 & $\ker\mathbf{A}^\top\subseteq\mathbb{R}^m$ & $m-r$ & $\{\mathbf{u}_{r+1},\dots,\mathbf{u}_m\}$ \\
\bottomrule
\end{tabular}
\end{table}

\noindent 几何骨架：把输入正交分解 $\mathbf{x}=\mathbf{x}_{\mathrm{row}}+\mathbf{x}_{\mathrm{null}}$，则 $\mathbf{A}\mathbf{x}=\mathbf{A}\mathbf{x}_{\mathrm{row}}$——\emph{零空间分量被丢弃，行空间被一一映到列空间}（$\mathbf{v}_i\mapsto\sigma_i\mathbf{u}_i$，可逆）。这把抽象的秩--零度变成可视的子空间图。

\begin{insight}[奇异值 vs 特征值：分析“形变”用前者，“迭代”用后者]
对非正规矩阵，存在向量被放大的倍数\emph{远超}最大特征值的模：$\mathbf{A}=\bigl[\begin{smallmatrix}1&100\\0&1\end{smallmatrix}\bigr]$ 特征值都是 1，但 $\mathbf{A}\bigl[\begin{smallmatrix}0\\1\end{smallmatrix}\bigr]=\bigl[\begin{smallmatrix}100\\1\end{smallmatrix}\bigr]$ 暴涨约 100 倍，最大奇异值 $\sigma_1\approx100$ 如实捕捉。所以衡量矩阵“放大能力 / 病态程度”必须用奇异值——条件数 $\kappa=\sigma_{\max}/\sigma_{\min}$ 据此定义。仅当 $\mathbf{A}$ 对称半正定时 SVD 与谱分解重合（$\mathbf{U}=\mathbf{V}=\mathbf{Q}$，$\boldsymbol{\Sigma}=\boldsymbol{\Lambda}$）；对称不定时奇异值 $=|\lambda_i|$，符号被吸收进 $\mathbf{U}\neq\mathbf{V}$。机器人学绝大多数是“一次性映射”（坐标变换、速度映射、投影），故 SVD 远比特征分解常用。
\end{insight}

\paragraph{极分解：把映射分成纯旋转 + 纯拉伸。}
SVD 三段把伸缩夹在两个\emph{无关}旋转中间；极分解把它重组成更对称的“一次旋转 $+$ 一次对称拉伸”。

\begin{theorem}[极分解]\label{thm:polar}
方阵 $\mathbf{A}\in\mathbb{R}^{n\times n}$ 可分解为 $\mathbf{A}=\mathbf{Q}\mathbf{P}$，$\mathbf{Q}$ 正交、$\mathbf{P}$ 对称半正定。$\mathbf{P}=\sqrt{\mathbf{A}^\top\mathbf{A}}$ \emph{永远唯一}；$\mathbf{Q}$ 唯一 $\iff$ $\mathbf{A}$ 满秩。由 SVD 三行导出：
\begin{equation}
  \mathbf{A}=\mathbf{U}\boldsymbol{\Sigma}\mathbf{V}^\top=\underbrace{(\mathbf{U}\mathbf{V}^\top)}_{\mathbf{Q}}\underbrace{(\mathbf{V}\boldsymbol{\Sigma}\mathbf{V}^\top)}_{\mathbf{P}}.
  \label{eq:polar}
\end{equation}
\end{theorem}

\begin{theorem}[极分解正交因子的最优性]\label{thm:polar-optimal}
在 Frobenius 范数下，$\mathbf{Q}=\mathbf{U}\mathbf{V}^\top$ 是离 $\mathbf{A}$ \emph{最近}的正交矩阵：$\mathbf{Q}=\arg\min_{\boldsymbol{\Omega}^\top\boldsymbol{\Omega}=\mathbf{I}}\|\mathbf{A}-\boldsymbol{\Omega}\|_F$。
\end{theorem}
\begin{proof}
$\|\mathbf{A}-\boldsymbol{\Omega}\|_F^2=\tr(\mathbf{A}^\top\mathbf{A})-2\tr(\mathbf{A}^\top\boldsymbol{\Omega})+\tr(\mathbf{I})$，前后两项与 $\boldsymbol{\Omega}$ 无关，故 $\min\|\mathbf{A}-\boldsymbol{\Omega}\|_F\iff\max\tr(\mathbf{A}^\top\boldsymbol{\Omega})$。代入 $\mathbf{A}=\mathbf{U}\boldsymbol{\Sigma}\mathbf{V}^\top$，记 $\mathbf{Z}=\mathbf{U}^\top\boldsymbol{\Omega}\mathbf{V}$（正交），$\tr(\mathbf{A}^\top\boldsymbol{\Omega})=\tr(\boldsymbol{\Sigma}\mathbf{Z})=\sum_i\sigma_iZ_{ii}\le\sum_i\sigma_i$（正交矩阵 $|Z_{ii}|\le1$），等号 $\iff\mathbf{Z}=\mathbf{I}\iff\boldsymbol{\Omega}=\mathbf{U}\mathbf{V}^\top=\mathbf{Q}$。
\end{proof}

\begin{pitfall}[陷阱：清洗旋转矩阵不能用归一化或 Gram--Schmidt]
IMU 积分、姿态估计中旋转矩阵 $\mathbf{R}$ 经大量浮点运算会失去正交性（$\mathbf{R}^\top\mathbf{R}\neq\mathbf{I}$）。\textbf{逐列归一化}只保证每列长度为 1、\emph{不保证列间正交}；\textbf{Gram--Schmidt} 虽得正交矩阵但\emph{偏心}（过度信任第一列、依赖列序），不是离 $\mathbf{R}$ 最近的。\textbf{正确做法}：对脏 $\mathbf{R}$ 做极分解（或 SVD 后令 $\boldsymbol{\Sigma}\to\mathbf{I}$）取 $\mathbf{Q}=\mathbf{U}\mathbf{V}^\top$（\cref{thm:polar-optimal} 保证它各向公平地最近）。\emph{但}极分解只保证 $\mathbf{Q}\in O(n)$，$\det\mathbf{Q}=\pm1$；$\det\mathbf{A}<0$ 时 $\mathbf{Q}$ 是\emph{反射}而非旋转，需额外 $\det$ 修正强制 $\det=+1$（见 \cref{thm:procrustes}）——否则机器人“镜像翻转”。
\end{pitfall}

\paragraph{Moore--Penrose 伪逆与最小二乘。}
SVD 直接给出\emph{伪逆}
\begin{equation}
  \mathbf{A}^+=\mathbf{V}\boldsymbol{\Sigma}^+\mathbf{U}^\top,\qquad\boldsymbol{\Sigma}^+=\diag(1/\sigma_1,\dots,1/\sigma_r,0,\dots,0).
  \label{eq:pinv}
\end{equation}
它把“行空间 $\leftrightarrow$ 列空间的可逆映射”反过来（$\mathbf{u}_i\mapsto\mathbf{v}_i/\sigma_i$），对两个零空间无能为力（留零）。$\mathbf{A}^+\mathbf{b}$ 统一了\emph{最小二乘}（超定时给残差最小解，与法方程 \cref{eq:normal-equation} 一致：列满秩时 $\mathbf{A}^+=(\mathbf{A}^\top\mathbf{A})^{-1}\mathbf{A}^\top$）与\emph{最小范数}（欠定时给范数最小解）。机械臂逆运动学 $\dot{\boldsymbol{\theta}}=\mathbf{J}^+\mathbf{V}_d$ 就用它。\emph{Eckart--Young}：截断 SVD $\sum_{i=1}^k\sigma_i\mathbf{u}_i\mathbf{v}_i^\top$ 是秩 $\le k$ 的\emph{最优}低秩逼近（PCA 的代数核心）。

\begin{theorem}[正交 Procrustes / Kabsch--Arun]\label{thm:procrustes}
给定两组对应点（去质心后）$\{\mathbf{x}_i\},\{\mathbf{y}_i\}$，求 $\mathbf{R}\in\SO(3)$ 最小化 $\sum_i\|\mathbf{y}_i-\mathbf{R}\mathbf{x}_i\|^2$。设互协方差 $\mathbf{H}=\sum_i\mathbf{y}_i\mathbf{x}_i^\top=\mathbf{U}\boldsymbol{\Sigma}\mathbf{V}^\top$，则最优
\begin{equation}
  \mathbf{R}^\star=\mathbf{U}\diag\!\bigl(1,1,\det(\mathbf{U}\mathbf{V}^\top)\bigr)\mathbf{V}^\top.
  \label{eq:procrustes}
\end{equation}
\end{theorem}
\begin{proof}[证明骨架]
$\sum_i\|\mathbf{y}_i-\mathbf{R}\mathbf{x}_i\|^2=\mathrm{const}-2\tr(\mathbf{R}\mathbf{H}^\top)$ 故 $\min\iff\max\tr(\mathbf{R}\mathbf{H}^\top)=\max\tr(\boldsymbol{\Sigma}\mathbf{U}^\top\mathbf{R}\mathbf{V})$，与 \cref{thm:polar-optimal} 同型，无约束最优为 $\mathbf{U}\mathbf{V}^\top$。但若 $\det(\mathbf{U}\mathbf{V}^\top)=-1$（反射），需在 $\SO(3)$ 内取次优——把最小奇异值对应方向反号，即得 \cref{eq:procrustes} 的 $\diag(1,1,\det(\mathbf{U}\mathbf{V}^\top))$ 修正项。
\end{proof}

\begin{insight}[那个神秘的 $\det$ 修正项为何必须存在]
SVD 的 $\mathbf{U},\mathbf{V}$ 列各有符号自由度，$\mathbf{U}\mathbf{V}^\top$ 可能落在 $O(3)$ 的\emph{反射}支（$\det=-1$）而非旋转支。物理上旋转必须 $\det=+1$；\cref{eq:procrustes} 的 $\diag(1,1,\det(\mathbf{U}\mathbf{V}^\top))$ 在反射情形把代价最小的那个方向翻回来，强制落入 $\SO(3)$。这正是 ICP/点云配准里“机器人突然镜像翻转”bug 的根源与解药——理解 SVD 的不唯一性是理解这个修正项的前提。同样的 $\det$ 修正也用于 \cref{thm:polar-optimal} 把极分解的正交因子约束到 $\SO(3)$。
\end{insight}

\begin{example}[可操作度椭球与奇异位形]\label{ex:manipulability}
机械臂雅可比 $\mathbf{J}=\mathbf{U}\boldsymbol{\Sigma}\mathbf{V}^\top$ 把单位关节速度球映成任务空间的\emph{可操作度椭球}：椭球长轴（大 $\sigma$）是末端“最易快速移动”方向，短轴（小 $\sigma$）是“最迟钝”方向。Yoshikawa 可操作度 $w=\sqrt{\det(\mathbf{J}\mathbf{J}^\top)}=\prod_i\sigma_i$ 是椭球体积。奇异位形判据 $\sigma_{\min}\to0$：椭球塌陷、某方向失去运动能力。冗余机械臂零空间运动方向就是零奇异值对应的右奇异向量。\textbf{注意椭球画在任务空间（输出），轴用 $\mathbf{U}$ 的列}；用错成 $\mathbf{V}$ 会把椭球画到关节空间、物理意义全错。
\end{example}

\begin{practice}
\begin{enumerate}
  \item 对 $\mathbf{A}=\bigl[\begin{smallmatrix}1&1\\0&1\\1&0\end{smallmatrix}\bigr]$ 手算完整 SVD：算 $\mathbf{A}^\top\mathbf{A}$ 的特征值/特征向量得 $\sigma_i,\mathbf{v}_i$；用 $\mathbf{u}_i=\mathbf{A}\mathbf{v}_i/\sigma_i$ 算 $\mathbf{u}_i$；扩成 $\mathbb{R}^3$ 正交基；验证 $\mathbf{U}\boldsymbol{\Sigma}\mathbf{V}^\top=\mathbf{A}$。
  \item 证明 $\mathbf{A}^\top\mathbf{A}$ 与 $\mathbf{A}\mathbf{A}^\top$ 有相同非零特征值（提示：$\mathbf{A}^\top\mathbf{A}\mathbf{v}=\lambda\mathbf{v}$ 两边左乘 $\mathbf{A}$）。据此说明算瘦高矩阵 $480\times3$ 的奇异值应算小的 $\mathbf{A}^\top\mathbf{A}$。
  \item 验证列满秩时 $\mathbf{A}^+=(\mathbf{A}^\top\mathbf{A})^{-1}\mathbf{A}^\top$，与法方程 \cref{eq:normal-equation} 解一致；并说明 $\mathbf{A}\mathbf{A}^+$ 是到 $\im\mathbf{A}$ 的正交投影。
\end{enumerate}
\end{practice}

% ════════════════════════════════════════════════════════════════
\section{极小多项式、Cayley--Hamilton 与 Jordan 标准形}
\label{sec:jordan}

\paragraph{动机：一切矩阵指数闭式的代数总开关。}
\cref{ch:lie} 的罗德里格斯公式 $\mathbf{R}=\mathbf{I}+\sin\theta\,\mathbf{a}^\wedge+(1-\cos\theta)(\mathbf{a}^\wedge)^2$ 只有三项、$\SE(3)$ 指数也是有限项闭式——为什么无穷级数 $\exp(\mathbf{M})=\sum\mathbf{M}^k/k!$ 会“坍缩”？答案在本节：\emph{极小多项式}把矩阵的高次幂归约为低次幂的线性组合，于是任何 $\mathbf{M}$ 的幂级数都坍缩成有限项。本节还给出一般（非正规、不可对角化）算子的结构分类——Jordan 标准形，它决定线性系统 $\dot{\mathbf{x}}=\mathbf{A}\mathbf{x}$ 的稳定性。

\paragraph{零化理想与极小多项式。}
对 $\mathbf{A}\in\mathbb{F}^{n\times n}$，$\{\mathbf{I},\mathbf{A},\mathbf{A}^2,\dots\}$ 在 $n^2$ 维的 $M_n(\mathbb{F})$ 中必线性相关，故存在非零多项式 $p$ 使 $p(\mathbf{A})=\mathbf{0}$。这样的 $p$ 构成\emph{零化理想} $\mathrm{Ann}(\mathbf{A})\subset\mathbb{F}[x]$。由于 $\mathbb{F}[x]$ 是\emph{主理想整环}（每个理想由单个元素生成——带余除法保证），$\mathrm{Ann}(\mathbf{A})=(m_{\mathbf{A}})$ 由唯一的首一最低次\emph{极小多项式} $m_{\mathbf{A}}$ 生成。它满足 $m_{\mathbf{A}}\mid p_{\mathbf{A}}$（整除特征多项式 $p_{\mathbf{A}}(\lambda)=\det(\lambda\mathbf{I}-\mathbf{A})$），二者有相同根集（但重数可不同）；$\mathbf{A}$ 可对角化 $\iff$ $m_{\mathbf{A}}$ 在 $\mathbb{F}$ 上分裂为互异一次因子。\footnote{完整条件是"$m_{\mathbf{A}}$ 在 $\mathbb{F}$ 上分裂\emph{且}无重根"。当 $\mathbb{F}=\mathbb{C}$（或更一般地 $\mathbb{F}$ 代数闭）时分裂自动成立，故简化为"$m_{\mathbf{A}}$ 无重根"即可。但在 $\mathbb{F}=\mathbb{R}$ 上分裂性不可省：如 $90^\circ$ 旋转 $\bigl[\begin{smallmatrix}0&-1\\1&0\end{smallmatrix}\bigr]$ 的 $m_{\mathbf{A}}=x^2+1$ 虽无重根却不在 $\mathbb{R}$ 上分裂，故\emph{不}可在 $\mathbb{R}$ 上对角化（须复化方可）。}

\begin{theorem}[Cayley--Hamilton]\label{thm:cayley-hamilton}
任意方阵满足自身的特征方程：$p_{\mathbf{A}}(\mathbf{A})=\mathbf{0}$。
\end{theorem}
\begin{proof}[证明骨架（Schur 路径，$\mathbb{F}=\mathbb{C}$）]
由 Schur 三角化（\cref{sec:spectral}）$\mathbf{A}$ 相似于上三角 $\mathbf{T}$，对角元为特征值 $\lambda_1,\dots,\lambda_n$，$p_{\mathbf{A}}(\lambda)=\prod_i(\lambda-\lambda_i)$。在不变旗 $V_0\subset V_1\subset\cdots\subset V_n=V$（$\dim V_k=k$，$\mathbf{A}V_k\subseteq V_k$）上，$(\mathbf{A}-\lambda_k\mathbf{I})V_k\subseteq V_{k-1}$（上三角性）。故复合 $(\mathbf{A}-\lambda_1\mathbf{I})\cdots(\mathbf{A}-\lambda_n\mathbf{I})$ 把 $V=V_n$ 逐级降到 $V_0=\{\mathbf{0}\}$，因各因子可交换（都是 $\mathbf{A}$ 的多项式）即得 $p_{\mathbf{A}}(\mathbf{A})=\mathbf{0}$。一般域用伴随矩阵（“望远镜求和”）路径。
\end{proof}

\begin{insight}[最小多项式是闭式收敛的“总开关”]
设 $\mathbf{M}$ 的极小多项式次数 $=d$，则 $\mathbf{M}^d$ 可由 $\{\mathbf{I},\mathbf{M},\dots,\mathbf{M}^{d-1}\}$ 线性表出，于是\emph{任何}关于 $\mathbf{M}$ 的幂级数（$\exp,\cos,(\mathbf{I}-\mathbf{M})^{-1}$、左雅可比 $\mathbf{J}$、伴随 $\mathrm{Ad}$）都坍缩成这 $d$ 个基的有限组合，系数是把级数按基归并后的标量函数。这就是为什么 $\so(3)$ 的一切闭式都落在 $\{\mathbf{I},\boldsymbol{\phi}^\wedge,(\boldsymbol{\phi}^\wedge)^2\}$ 三族（$d=3$）。遇到新李群（如 $\mathrm{SE}_2(3)$）也能照方抓药：先求特征值定极小多项式次数，闭式的“基”和项数当即可知（对照 \cref{ch:lie} 的最小多项式小节）。
\end{insight}

\begin{theorem}[Jordan 标准形]\label{thm:jordan}
设 $\mathbb{F}$ 代数闭（或 $p_{\mathbf{A}}$ 在 $\mathbb{F}$ 上分裂）。则 $\mathbf{A}$ 相似于 Jordan 块的直和 $\mathbf{J}=\bigoplus_{i}\mathbf{J}_{k_i}(\lambda_i)$，其中 $\mathbf{J}_k(\lambda)=\lambda\mathbf{I}_k+\mathbf{N}_k$（$\mathbf{N}_k$ 超对角线为 1 的幂零块）。块的尺寸由 $\dim\ker(\mathbf{A}-\lambda\mathbf{I})^j$ 的递推唯一确定。
\end{theorem}

\noindent 证明路线：\emph{准素分解} $V=\bigoplus_\lambda\ker(\mathbf{A}-\lambda\mathbf{I})^{a_\lambda}$（按特征值拆成广义特征空间，由 $m_{\mathbf{A}}=\prod p_i^{m_i}$ 与 Bézout 给出投影）$\to$ 每个广义特征空间上 $\mathbf{A}=\lambda\mathbf{I}+\mathbf{N}$（$\mathbf{N}$ 幂零）$\to$ 对幂零部分用核旗升法构造 Jordan 链。\emph{有理标准形}（不依赖代数闭域）由 $\mathbb{F}[x]$-模在 PID 上的结构定理给出不变因子分解 $V\cong\bigoplus\mathbb{F}[x]/(d_i)$，把 Jordan、Cayley--Hamilton、$m_{\mathbf{A}}\mid p_{\mathbf{A}}$ 统一为同一代数定理的不同面目。Jordan 基一般\emph{不正交}（要正交性须用 Schur 三角化）。

\paragraph{矩阵指数与罗德里格斯公式。}
$\dot{\mathbf{x}}=\mathbf{A}\mathbf{x}$ 的解 $\mathbf{x}(t)=e^{\mathbf{A}t}\mathbf{x}(0)$。借 Jordan 形 $e^{\mathbf{A}t}=\mathbf{S}e^{\mathbf{J}t}\mathbf{S}^{-1}$，每个 $k\times k$ Jordan 块给 $e^{\mathbf{J}_k(\lambda)t}=e^{\lambda t}\sum_{j=0}^{k-1}\frac{t^j}{j!}\mathbf{N}_k^j$（$\mathbf{N}_k$ 幂零截断）。稳定性由 Jordan 结构决定：$\dot{\mathbf{x}}=\mathbf{A}\mathbf{x}$ 渐近稳定 $\iff$ 全特征值实部 $<0$；$\lambda=0$ 的 $2\times2$ 块产生 $t$ 增长项导致漂移（仅看特征值不够）。

\begin{theorem}[罗德里格斯作为 $\so(3)$ 上的 Cayley--Hamilton]\label{thm:rodrigues-alg}
单位轴 $\mathbf{a}$（$\|\mathbf{a}\|=1$）的反对称矩阵 $\mathbf{a}^\wedge\in\so(3)$ 满足极小多项式 $(\mathbf{a}^\wedge)^3=-\mathbf{a}^\wedge$。代入指数级数并按奇偶项归并，给出闭式
\begin{equation}
  e^{\theta\,\mathbf{a}^\wedge}=\mathbf{I}+\sin\theta\,\mathbf{a}^\wedge+(1-\cos\theta)(\mathbf{a}^\wedge)^2.
  \label{eq:rodrigues-alg}
\end{equation}
\end{theorem}
\begin{proof}
由 $(\mathbf{a}^\wedge)^2=\mathbf{a}\mathbf{a}^\top-\mathbf{I}$（单位向量恒等式）得 $(\mathbf{a}^\wedge)^3=(\mathbf{a}^\wedge)(\mathbf{a}\mathbf{a}^\top-\mathbf{I})=-\mathbf{a}^\wedge$（因 $\mathbf{a}^\wedge\mathbf{a}=\mathbf{0}$）。于是 $\mathbf{a}^\wedge$ 的幂在 $\{\mathbf{a}^\wedge,(\mathbf{a}^\wedge)^2\}$ 间循环：奇次项凑成 $\sin\theta\,\mathbf{a}^\wedge$，偶次项凑成 $(1-\cos\theta)(\mathbf{a}^\wedge)^2$，常数项 $\mathbf{I}$，即 \cref{eq:rodrigues-alg}。它与 \cref{ch:rigid_body}、\cref{ch:lie} 的罗德里格斯公式一字不差。
\end{proof}

\begin{insight}[类比：罗德里格斯之于 $\so(3)$ 犹如欧拉公式之于复数]
\cref{eq:rodrigues-alg} 与欧拉公式 $e^{i\theta}=\cos\theta+i\sin\theta$ 是同一思想：都利用“极小多项式的循环性”（$i^2=-1$ 对应 $(\mathbf{a}^\wedge)^3=-\mathbf{a}^\wedge$）把指数截断为有限三角表达。边界：欧拉是标量等式、罗德里格斯是矩阵等式；$i^2=-1$ 是二次关系、$(\mathbf{a}^\wedge)^3=-\mathbf{a}^\wedge$ 是三次关系（故罗德里格斯有三项）。理解这个来源后，对\emph{任何}矩阵——只要知道其极小多项式——都能推导类似闭式，而非死记孤立公式。$\SE(3)$ 指数含 $\mathbf{G}(\theta)\boldsymbol{\rho}$ 项的闭式同理（\cref{ch:lie}）。
\end{insight}

\begin{practice}
\begin{enumerate}
  \item 对 $\bigl[\begin{smallmatrix}2&1&0\\0&2&1\\0&0&3\end{smallmatrix}\bigr]$ 求 Jordan 标准形与 Jordan 基；对 90$^\circ$ 旋转 $\bigl[\begin{smallmatrix}0&-1\\1&0\end{smallmatrix}\bigr]$ 说明为何在 $\mathbb{R}$ 上无 Jordan 形，写出复 Jordan 形。
  \item 由 $\mathbf{a}^\wedge$（$\|\mathbf{a}\|=1$）验证 $(\mathbf{a}^\wedge)^3=-\mathbf{a}^\wedge$，推导 \cref{eq:rodrigues-alg}（草稿纸完成）。
  \item 用 Cayley--Hamilton 把 $\mathbf{A}^n$（$n\times n$）表为 $\{\mathbf{I},\mathbf{A},\dots,\mathbf{A}^{n-1}\}$ 的线性组合，解释能控性矩阵 $[\mathbf{B},\mathbf{A}\mathbf{B},\dots,\mathbf{A}^{n-1}\mathbf{B}]$ 为何只到 $n-1$ 次。
\end{enumerate}
\end{practice}

% ════════════════════════════════════════════════════════════════
\section{张量积、外代数与 \texorpdfstring{$\so(3)\cong\Lambda^2\mathbb{R}^3$}{so(3)=Lambda2R3}}
\label{sec:tensor}

\paragraph{动机：从单线性跳到多线性。}
前几节把“单线性”代数做到极致。但机器人学的日用品大量是\emph{多线性}的——对每个输入槽分别线性、变量之间\emph{耦合}：内积 $\langle\mathbf{u},\mathbf{v}\rangle$、叉积 $\mathbf{u}\times\mathbf{v}$、行列式 $\det(\mathbf{v}_1,\dots,\mathbf{v}_n)$、惯性张量、动能 $\frac12\dot{\mathbf{q}}^\top\mathbf{M}\dot{\mathbf{q}}$。本节用\emph{张量积}给多线性映射统一的代数框架，用\emph{外代数}给行列式以坐标无关定义，并揭示 $\so(3)\cong\Lambda^2\mathbb{R}^3$——角速度“伪向量”的代数真身。

\begin{definition}[张量积的泛性质]\label{def:tensor}
向量空间 $V,W$ 的\emph{张量积}是一个空间 $V\otimes W$ 配双线性映射 $\otimes:V\times W\to V\otimes W$，使得\emph{任意}双线性映射 $B:V\times W\to U$ 唯一地分解为 $B=\tilde{B}\circ\otimes$（$\tilde{B}:V\otimes W\to U$ 线性）。它至同构唯一，$\dim(V\otimes W)=\dim V\cdot\dim W$，基为 $\{\mathbf{e}_i\otimes\mathbf{f}_j\}$。
\end{definition}

\noindent 泛性质把“双线性”\emph{线性化}：$V\otimes W$ 是承接一切双线性映射的“万能中转站”。有限维有自然同构 $V^*\otimes W\cong\Hom(V,W)$（线性映射就是 $(1,1)$ 型张量），$\Hom(V\otimes W,U)\cong\Hom(V,\Hom(W,U))$（Curry 化）。物理中协方差 $\boldsymbol{\Sigma}\in V\otimes V$、信息矩阵 $\boldsymbol{\Lambda}\in V^*\otimes V^*$ 的上下标位置正源于此。

\paragraph{外代数与楔积。}
对张量代数 $T(V)=\bigoplus_k V^{\otimes k}$ 商去 $\langle\mathbf{v}\otimes\mathbf{v}\rangle$ 得\emph{外代数} $\Lambda(V)$，乘法 $\wedge$ 满足\emph{交替性} $\mathbf{v}\wedge\mathbf{v}=0$（故 $\mathbf{u}\wedge\mathbf{v}=-\mathbf{v}\wedge\mathbf{u}$，分级反交换 $\alpha\wedge\beta=(-1)^{pq}\beta\wedge\alpha$）。$\Lambda^k(V)$ 维数 $\binom{n}{k}$，$\Lambda^n(V)$ 一维，$\dim\Lambda(V)=2^n$。关键事实：$\mathbf{v}_1\wedge\cdots\wedge\mathbf{v}_k\neq0\iff\{\mathbf{v}_i\}$ 线性无关。

\begin{theorem}[行列式即 $\Lambda^n$]\label{thm:det-exterior}
线性映射 $\mathbf{T}:V\to V$（$\dim V=n$）诱导 $\Lambda^n(\mathbf{T})$ 作用在一维空间 $\Lambda^n(V)$ 上、即某个标量乘法，该标量\emph{定义}为 $\det\mathbf{T}$：
\begin{equation}
  \mathbf{T}\mathbf{v}_1\wedge\cdots\wedge\mathbf{T}\mathbf{v}_n=(\det\mathbf{T})\,\mathbf{v}_1\wedge\cdots\wedge\mathbf{v}_n.
  \label{eq:det-exterior}
\end{equation}
\end{theorem}

\noindent 这个定义\emph{不依赖任何基}（$\mathbf{T}$ 与 $V$ 就够），$|\det\mathbf{T}|$ 是体积放大率、$\mathrm{sgn}(\det\mathbf{T})$ 是定向保持/反转。乘法性 $\det(\mathbf{A}\mathbf{B})=\det\mathbf{A}\det\mathbf{B}$ 由函子性\emph{一行}得证：$\Lambda^n(\mathbf{A}\mathbf{B})=\Lambda^n(\mathbf{A})\circ\Lambda^n(\mathbf{B})=(\det\mathbf{A})(\det\mathbf{B})\cdot\mathrm{id}$；取标准基展开 \cref{eq:det-exterior} 即恢复排列公式 $\det\mathbf{T}=\sum_\sigma\mathrm{sgn}(\sigma)\prod_i a_{\sigma(i),i}$。这解释了 $\SO(n)=\{\mathbf{R}\in O(n):\det\mathbf{R}=+1\}$——正行列式即保持定向。

\begin{theorem}[Hodge 星与 $\so(3)\cong\Lambda^2\mathbb{R}^3$]\label{thm:hodge}
配内积与定向后，\emph{Hodge 星} $\star:\Lambda^k(V)\to\Lambda^{n-k}(V)$ 由 $\alpha\wedge\star\beta=\langle\alpha,\beta\rangle\,\mathrm{vol}$ 唯一确定，是等距同构（利用 $\binom{n}{k}=\binom{n}{n-k}$）。在 $\mathbb{R}^3$ 上 $\star:\Lambda^2\mathbb{R}^3\xrightarrow{\sim}\Lambda^1\mathbb{R}^3=\mathbb{R}^3$，给出同构链
\begin{equation}
  \mathbb{R}^3\xrightarrow{\ \wedge\ }\Lambda^2\mathbb{R}^3\xrightarrow{\ \star\ }\mathbb{R}^3,\qquad
  \mathbb{R}^3\cong\Lambda^2\mathbb{R}^3\cong\so(3),
  \label{eq:so3-iso}
\end{equation}
且叉积 $\mathbf{u}\times\mathbf{v}=\star(\mathbf{u}\wedge\mathbf{v})$、帽映射 $\boldsymbol{\omega}^\wedge\mathbf{v}=\boldsymbol{\omega}\times\mathbf{v}$ 恰是这条链。
\end{theorem}

\begin{insight}[叉积是三维的“偶然”，角速度其实是 2-形式]
$\dim\Lambda^2\mathbb{R}^n=\binom{n}{2}$ 仅在 $n=3$ 时等于 $n$，这\emph{唯一}地使“把两个向量变成一个向量”的叉积存在；$\mathbb{R}^4$ 中 $\dim\Lambda^2\mathbb{R}^4=6\neq4$，\emph{没有}叉积（六个独立旋转平面对应 $\so(4)$ 的六维）。所以叉积不是基本运算，而是楔积（纯代数、坐标无关）后接 Hodge 星（依赖度量与定向）的复合；$\mathbf{u}\wedge\mathbf{v}\in\Lambda^2\mathbb{R}^3$ 才是“真对象”（带定向的面积元）。角速度 $\boldsymbol{\omega}$ 实为 2-形式的 Hodge 对偶——这正是 \cref{eq:so3-iso} 揭示的：$\boldsymbol{\omega}\in\mathbb{R}^3$ 与反对称 $\boldsymbol{\omega}^\wedge\in\so(3)$ 的换算 $(\cdot)^\wedge$ 就是 $\mathbb{R}^3\cong\Lambda^2\mathbb{R}^3\cong\so(3)$ 的同构链。这为 \cref{ch:lie} 的反对称化记号、螺旋 $\mathbf{S}\in\se(3)\cong\mathbb{R}^3\oplus\Lambda^2\mathbb{R}^3$、视觉 SLAM 线特征的 Plücker 坐标（$\Lambda^2\mathbb{R}^4$）提供统一的代数根基。$\SO(4)/\SE(3)$ 中此“简便”失效（$\dim\so(4)=6\neq4$），改用 Plücker 坐标或伴随 $\mathrm{ad}_{\boldsymbol{\xi}}$。
\end{insight}

\begin{example}[惯性张量与 Plücker 坐标的代数身份]\label{ex:tensor-robot}
惯性张量 $I_{ij}=\int\rho(r^2\delta_{ij}-r_ir_j)\,dV$ 是 $V\otimes V$ 中的\emph{对称} 2-张量（对称代数 $S^2(V)$ 元素），其变换律由张量积函子性给出。Plücker 线坐标是 $\Lambda^2\mathbb{R}^4$ 中的单纯 2-向量（满足 Plücker 关系），用于视觉 SLAM 的线特征。空间惯性矩阵 $\mathbf{I}\in\mathbb{R}^{6\times6}$（Featherstone 空间代数）对称正定，可谱分解出主惯性轴。这些“张量”不再是模糊的物理符号，而是 \cref{def:tensor} 与外代数中的严格代数对象。
\end{example}

\begin{practice}
\begin{enumerate}
  \item 用外代数定义证明 $\det(c\mathbf{A})=c^n\det\mathbf{A}$ 与“$\mathbf{A}$ 可逆 $\iff\det\mathbf{A}\neq0$”。
  \item 在 $\mathbb{R}^3$ 标准基下验证 $\star(\mathbf{e}_1\wedge\mathbf{e}_2)=\mathbf{e}_3$ 等，并据此从 $\mathbf{u}\times\mathbf{v}=\star(\mathbf{u}\wedge\mathbf{v})$ 推出叉积分量公式。
  \item 证明 $n$ 个向量的 Gram 行列式 $\det(\langle\mathbf{v}_i,\mathbf{v}_j\rangle)=\|\mathbf{v}_1\wedge\cdots\wedge\mathbf{v}_n\|^2$，并联系可操作度 $w=\sqrt{\det(\mathbf{J}\mathbf{J}^\top)}$。
\end{enumerate}
\end{practice}

% ════════════════════════════════════════════════════════════════
\section*{本章常见误解汇总}
\begin{table}[H]\centering\small
\begin{tabular}{p{0.46\textwidth} p{0.46\textwidth}}
\toprule
\textbf{常见误解} & \textbf{正确理解} \\
\midrule
$V^*\cong V$ 所以二者没区别 & 同构需选基（非自然）；元素物理含义与变换律不同（\cref{def:dual}）\\
$\mathbf{J}^\top$ 是内积伴随，依赖度量 & 是对偶映射，由虚功原理给出，\emph{不依赖}度量（\cref{ex:jacobian-transpose}）\\
对偶 $\mathbf{T}^\top$ 与伴随 $\mathbf{A}^*$ 是一回事 & 前者无需内积，后者需要；仅实正交基下矩阵一致 \\
零化子 $W^\circ$ 等于正交补 $W^\perp$ & $W^\circ\subseteq V^*$、$W^\perp\subseteq V$；后者需内积（\cref{eq:annihilator-duality}）\\
互易螺旋系就是欧氏正交补 & 是做功（Klein）配对下的零化子，与互易系维数和恒为 $6$，不需内积（\cref{thm:recip-dim}）\\
$\mathbf{R}^{-1}=\mathbf{R}^\top$ 是巧合 & 正交即保内积的必然（\cref{ex:orthogonal}）\\
能对角化 = 能正交对角化 & 后者是对称矩阵专属（谱定理）；非对称即使可对角化也一般不正交 \\
特征值都是实数 & 仅对称（自伴）矩阵保证；旋转矩阵特征值含 $e^{\pm i\theta}$ \\
矩阵平方根 = 逐元素开方 & 必须谱分解、对特征值开方（\cref{cor:sqrt}）\\
奇异值 = 特征值的绝对值 & 仅正规矩阵成立；一般天差地别（\cref{eq:svd} 后注）\\
SVD 唯一 & 奇异值唯一，$\mathbf{U},\mathbf{V}$ 不唯一（符号 + 重值子空间旋转）\\
清洗旋转用归一化 / Gram--Schmidt & 应用极分解取 $\mathbf{Q}=\mathbf{U}\mathbf{V}^\top$（最近正交，\cref{thm:polar-optimal}）\\
极分解的 $\mathbf{Q}$ 一定是旋转 & 只保证 $\in O(n)$；$\det\mathbf{A}<0$ 时是反射，需 $\det$ 修正（\cref{thm:procrustes}）\\
$(\mathbf{a}^\wedge)^3=-\mathbf{a}^\wedge$ 是孤立技巧 & 是 $\so(3)$ 极小多项式，闭式收敛的统一源头（\cref{thm:rodrigues-alg}）\\
线性系统稳定只看特征值 & 重特征值需看 Jordan 块大小（$\lambda=0$ 的 $2\times2$ 块致漂移）\\
叉积在任意维都存在 & 仅 $n=3$（与 $n=7$）；因 $\binom{3}{2}=3$（\cref{thm:hodge}）\\
行列式只是计算工具 & 是 $\Lambda^n$ 上的标量，体积与定向的合体（\cref{thm:det-exterior}）\\
\bottomrule
\end{tabular}
\end{table}

% ════════════════════════════════════════════════════════════════
\section*{本章小结}

\begin{table}[H]\centering\small
\caption{符号表}
\begin{tabular}{lll}
\toprule
符号 & 含义 & 首次出现 \\
\midrule
$\mathbb{F},\ V,W,U$ & 域 / 向量空间 & \cref{def:vector-space} \\
$W\le V,\ V/W,\ \oplus$ & 子空间 / 商空间 / 直和 & \cref{sec:vector-space} \\
$\ker\mathbf{T},\ \im\mathbf{T},\ \mathrm{rank}$ & 核 / 像 / 秩 & \cref{thm:rank-nullity} \\
$V^*,\ \langle f,\mathbf{v}\rangle,\ \mathbf{T}^\top,\ W^\circ$ & 对偶空间 / 配对 / 对偶映射 / 零化子 & \cref{def:dual} \\
$\odot,\ \boldsymbol{\Pi},\ \mathcal{S}^{\perp r}$ & 互易积 / Klein 型 / 互易螺旋系 & \cref{def:recip-product} \\
$\langle\cdot,\cdot\rangle,\ \|\cdot\|,\ W^\perp$ & 内积 / 范数 / 正交补 & \cref{def:inner-product} \\
$\mathbf{A}^*,\ \proj_W$ & 伴随算子 / 正交投影 & \cref{eq:adjoint},\cref{thm:projection} \\
$\mathbf{Q}\boldsymbol{\Lambda}\mathbf{Q}^\top$ & 实对称谱分解 & \cref{eq:spectral} \\
$\mathbf{A}^{1/2},\ \mathbf{L}\mathbf{L}^\top$ & 对称半正定平方根 / Cholesky & \cref{cor:sqrt} \\
$\mathbf{U}\boldsymbol{\Sigma}\mathbf{V}^\top,\ \sigma_i$ & SVD / 奇异值 & \cref{eq:svd} \\
$\mathbf{A}=\mathbf{Q}\mathbf{P},\ \mathbf{A}^+$ & 极分解 / Moore--Penrose 伪逆 & \cref{eq:polar},\cref{eq:pinv} \\
$m_{\mathbf{A}},\ p_{\mathbf{A}},\ \mathbf{J}_k(\lambda)$ & 极小 / 特征多项式 / Jordan 块 & \cref{sec:jordan} \\
$V\otimes W,\ \Lambda^k(V),\ \wedge,\ \star$ & 张量积 / 外幂 / 楔积 / Hodge 星 & \cref{def:tensor},\cref{thm:hodge} \\
$\so(3),\ (\cdot)^\wedge,\ \times$ & 旋转李代数 / 反对称化 / 叉积 & \cref{eq:so3-iso} \\
\bottomrule
\end{tabular}
\end{table}

\begin{table}[H]\centering\small
\caption{定理速查表}
\begin{tabular}{lll}
\toprule
定理 / 公式 & 一句话说明 & 对应 \\
\midrule
秩--零度 \cref{thm:rank-nullity} & 冗余度 $+$ 任务维数 $=$ 关节数 & $\dim\ker+\dim\im=\dim V$ \\
对偶映射 \cref{eq:dual-map} & $\boldsymbol{\tau}=\mathbf{J}^\top\mathbf{F}$，不需内积 & \cref{ex:jacobian-transpose} \\
行秩=列秩 \cref{eq:annihilator-duality} & 零化子对偶给概念证明 & $\mathrm{rank}\,\mathbf{T}=\mathrm{rank}\,\mathbf{T}^\top$ \\
互易螺旋系 \cref{thm:recip-dim} & $\dim\mathcal{S}+\dim\mathcal{S}^{\perp r}=6$，零化子推论 & 约束 $\leftrightarrow$ 自由度 \\
Riesz / 伴随 \cref{eq:adjoint} & 对偶与内积焊接，定义 $\mathbf{A}^*$ & \cref{thm:riesz} \\
法方程 \cref{eq:normal-equation} & 最小二乘 = 正交投影到列空间 & \cref{thm:projection} \\
实对称谱定理 \cref{eq:spectral} & 转主轴$\to$缩放$\to$转回，椭球摆正 & \cref{thm:spectral-real} \\
正定判据 \cref{thm:positive-definite} & 正定 $\iff$ 特征值全正（旋转不变） & \\
对称平方根 \cref{cor:sqrt} & 谱分解开方；Cholesky $\mathbf{L}\mathbf{L}^\top$ & 白化 / 采样高斯 \\
SVD \cref{eq:svd} & 任意映射 = 旋转$\times$伸缩$\times$旋转 & \cref{thm:svd} \\
极分解 \cref{eq:polar} & 旋转 $+$ 对称拉伸；$\mathbf{Q}$ 最近正交 & \cref{thm:polar-optimal} \\
伪逆 \cref{eq:pinv} & 统一最小二乘与最小范数 & $\dot{\boldsymbol{\theta}}=\mathbf{J}^+\mathbf{V}_d$ \\
Procrustes \cref{eq:procrustes} & ICP 最优旋转 $+\det$ 修正防反射 & \cref{thm:procrustes} \\
Cayley--Hamilton \cref{thm:cayley-hamilton} & $p_{\mathbf{A}}(\mathbf{A})=\mathbf{0}$，高次幂降阶 & \\
极小多项式 \cref{thm:rodrigues-alg} & 一切闭式收敛的总开关 & 罗德里格斯三项 \\
$\det=\Lambda^n$ \cref{eq:det-exterior} & 坐标无关、乘法性一行证 & \cref{thm:det-exterior} \\
$\so(3)\cong\Lambda^2\mathbb{R}^3$ \cref{eq:so3-iso} & 角速度是 2-形式；叉积是三维偶然 & \cref{thm:hodge} \\
\bottomrule
\end{tabular}
\end{table}

\begin{table}[H]\centering\small
\caption{知识点总表}
\begin{tabular}{clll}
\toprule
\# & 知识点 & 核心要点 & 难度 \\
\midrule
1 & 向量空间 / 商空间 & 公理化、秩--零度、相似不变量 & \(\bigstar\bigstar\) \\
2 & 对偶空间 & twist/wrench 对偶、$\mathbf{T}^\top$、零化子、互易螺旋系 & \(\bigstar\bigstar\bigstar\) \\
3 & 内积 / 正交投影 & 法方程 = 投影；QR；Cauchy--Schwarz & \(\bigstar\bigstar\bigstar\) \\
4 & 伴随算子 & Riesz；$\mathbf{A}^*$ vs $\mathbf{T}^\top$ vs $\mathrm{Ad}$ & \(\bigstar\bigstar\bigstar\bigstar\) \\
5 & 谱定理 / 正定 & 正交对角化、椭球、特征值判据 & \(\bigstar\bigstar\bigstar\) \\
6 & 平方根 / Cholesky & 谱分解开方；白化、采样高斯 & \(\bigstar\bigstar\bigstar\) \\
7 & SVD & $\mathbf{A}^\top\mathbf{A}$ 谱分解；四子空间；可操作度 & \(\bigstar\bigstar\bigstar\bigstar\) \\
8 & 极分解 / 伪逆 & 最近正交；Procrustes；最小二乘 & \(\bigstar\bigstar\bigstar\bigstar\) \\
9 & Cayley--Hamilton / Jordan & 极小多项式；矩阵指数；稳定性 & \(\bigstar\bigstar\bigstar\bigstar\bigstar\) \\
10 & 张量积 / 外代数 & $\det=\Lambda^n$；Hodge；$\so(3)\cong\Lambda^2\mathbb{R}^3$ & \(\bigstar\bigstar\bigstar\bigstar\bigstar\) \\
\bottomrule
\end{tabular}
\end{table}

% ════════════════════════════════════════════════════════════════
\section*{延伸阅读}
\begin{itemize}
  \item \textbf{现代算子论主线}：Axler《Linear Algebra Done Right》\cite{axler2024ladr}——向量空间、内积、谱定理、SVD 的无行列式现代讲法，几何直觉首选。
  \item \textbf{矩阵分析工具箱}：Horn \& Johnson《Matrix Analysis》\cite{hornjohnson2013}——SVD、扰动、范数、Jordan 的机器人日常工具箱；Strang《Introduction to Linear Algebra》\cite{strang2023intro}——四个基本子空间与五大分解的直觉串讲。
  \item \textbf{数值可靠性}：Trefethen \& Bau《Numerical Linear Algebra》\cite{trefethenbau1997}、Golub \& Van Loan《Matrix Computations》\cite{golubvanloan2013}——束调整 / 卡尔曼后端数值稳定性、SVD/QR/Cholesky 算法。
  \item \textbf{Jordan 与张量/外代数}：Hoffman \& Kunze《Linear Algebra》\cite{hoffmankunze1971}（有理 / Jordan 标准形金标准）、Greub《Multilinear Algebra》\cite{greub1978multilinear}（张量、外、Clifford 最深）、Halmos《Finite-Dimensional Vector Spaces》\cite{halmos1958fdvs}（对偶与谱的优雅入门）。
  \item \textbf{对接后续}：本章谱/SVD/伴随直接支撑 \cref{ch:lie}（罗德里格斯、伴随、最小多项式）、P1 凸优化与非线性最小二乘（法方程、Cholesky、KKT 对偶）；优化主线参考 Boyd \& Vandenberghe\cite{boyd2004convex}。
\end{itemize}

\section*{本章与后续章节的关系}
\begin{table}[H]\centering\small
\begin{tabular}{lll}
\toprule
后续章节 & 关系 & 本章铺垫 \\
\midrule
\cref{ch:lie} 李群与李代数 & 罗德里格斯、伴随、闭式收敛 & 极小多项式（\cref{thm:rodrigues-alg}）、$\so(3)\cong\Lambda^2\mathbb{R}^3$、谱/SVD \\
\cref{ch:rigid_body} 刚体运动 & 旋转矩阵、叉积、惯性张量 & 正交性（\cref{ex:orthogonal}）、张量/外代数（\cref{ex:tensor-robot}）\\
P1 凸优化 \cref{ch:convex} & KKT、对偶、二次型 & 对偶空间、正定（\cref{thm:positive-definite}）、投影（\cref{thm:projection}）\\
P1 非线性最小二乘 & 法方程、信息矩阵、Cholesky & \cref{eq:normal-equation}、对称平方根（\cref{cor:sqrt}）、伪逆 \\
概率与估计 \cref{ch:prob} & 协方差谱分解、Mahalanobis、白化 & 谱定理（\cref{ex:spectral-robot}）、SVD、$\boldsymbol{\Sigma}^{-1/2}$ \\
\bottomrule
\end{tabular}
\end{table}

\section*{故障排查手册}
\begin{enumerate}
  \item \textbf{症状}：旋转矩阵积分后\emph{逐渐失正交}（向量长度漂移）。\textbf{对策}：定期极分解取 $\mathbf{Q}=\mathbf{U}\mathbf{V}^\top$（\cref{thm:polar-optimal}），\emph{加 $\det$ 修正}强制 $\det=+1$（\cref{thm:procrustes}），勿用归一化 / Gram--Schmidt。
  \item \textbf{症状}：点云配准后机器人\emph{镜像翻转}。\textbf{原因}：$\det(\mathbf{U}\mathbf{V}^\top)=-1$ 落入反射支。\textbf{对策}：用 \cref{eq:procrustes} 的 $\diag(1,1,\det(\mathbf{U}\mathbf{V}^\top))$ 修正。
  \item \textbf{症状}：最小二乘 / 法方程\emph{病态或不收敛}。\textbf{排查}：检查条件数 $\kappa=\sigma_{\max}/\sigma_{\min}$（SVD）；病态时用 QR 或截断 SVD（伪逆 \cref{eq:pinv}）代替直接解 $\mathbf{A}^\top\mathbf{A}$；信息矩阵须对称正定方可 Cholesky（\cref{cor:sqrt}）。
  \item \textbf{症状}：“矩阵平方根 / 白化”结果\emph{不对称或平方不还原}。\textbf{原因}：误用逐元素开方。\textbf{对策}：谱分解后对\emph{特征值}开方（\cref{cor:sqrt}）。
  \item \textbf{症状}：线性系统特征值实部 $<0$ 却仍\emph{发散 / 漂移}。\textbf{原因}：$\lambda=0$ 或重特征值的非平凡 Jordan 块产生 $t$ 增长项。\textbf{对策}：查 Jordan 块大小，而非只看特征值（\cref{thm:jordan}）。
\end{enumerate}
